\chapter{Fourier transforms on \texorpdfstring{$L_2(-\infty,\infty)$}{L2(-infinity, infinity)}}
\label{ch:II}

\setcounter{section}{0}
\section{Introduction}
\label{sec:II-1}

The Banach space $L_{2}(-\infty,\infty)$ is endowed with an inner product. For any two functions $f,g$ belonging to $L_{2}(-\infty,\infty)$, the inner product $(f,g)$ is defined by
\begin{equation}\tag{1.1}\label{eq:II-1.1}
(f,g)=\int_{-\infty}^{\infty}f(x)\,\overline{g(x)}\,dx,
\end{equation}
the integral existing because of Schwarz's inequality. (The bar ``--'' denotes complex conjugation.) It has the following properties:
\begin{align*}
&\text{(i)}\quad \overline{(f,g)}=(g,f), && \text{(ii)}\quad (f,f)=\|f\|_{2}^{2},\\
&\text{(iii)}\quad (f,ag)=\overline{a}\,(f,g),\quad a\in\C, && \text{(iv)}\quad (f,g_{1}+g_{2})=(f,g_{1})+(f,g_{2}),\\
&\text{(v)}\quad |(f,g)|\le \|f\|_{2}\,\|g\|_{2}, && \text{(vi)}\quad (f,g)\text{ is continuous in }f,\text{ and in }g,
\end{align*}
for $f,g,g_{1},g_{2}\in L_{2}(-\infty,\infty)$.

A Banach space with such an inner product is referred to as a Hilbert space, and the particular Hilbert space $L_{2}(-\infty,\infty)$ is separable, in the sense that it contains a countable subset which is dense in $L_{2}(-\infty,\infty)$.

If a function $f$ belongs to $L_{2}(-\infty,\infty)$, it does not necessarily follow that $f\in L_{1}(-\infty,\infty)$, as for example $f(x)=\frac{1}{1+|x|}$. One has therefore to define the Fourier transform on $L_{2}(-\infty,\infty)$ somewhat differently from the way it was done in the case of $L_{1}(-\infty,\infty)$ in Chapter~I, but make sure that the two definitions coincide if $f\in L_{1}(-\infty,\infty)\cap L_{2}(-\infty,\infty)$.

We shall define a Fourier transform first on a dense subset of $L_{2}(-\infty,\infty)$, and then extend it uniquely to the whole of $L_{2}(-\infty,\infty)$. The extension will be defined almost everywhere, instead of everywhere as in the case of $L_{1}(-\infty,\infty)$. One can start with the dense subset of bounded functions belonging to $L_{1}(-\infty,\infty)$, or the subset $L_{1}(-\infty,\infty)\cap L_{2}(-\infty,\infty)$, or the subset of continuous functions which are of bounded variation over a finite interval and vanish outside that interval, or the subset of step-functions, or the subset $\mathcal S$ of infinitely differentiable functions which ``decrease rapidly'' (Schwartz's space, \S\ref{sec:I-3}, Ch.~I). For convenience we shall choose the last alternative.

The Fourier transform will turn out to be a bounded, linear operator on $\mathcal S$, which is an isometry, and which can be uniquely extended to a bounded linear operator on all of $L_{2}(-\infty,\infty)$, whose range coincides with $L_{2}(-\infty,\infty)$, and which is ``isometric'', hence a unitary operator. Thus if $f\in L_{2}(-\infty,\infty)$, its Fourier transform $F[f]$ automatically belongs to $L_{2}(-\infty,\infty)$. Because of the symmetry typified by this fact, we define $F[f]=\frac{1}{\sqrt{2\pi}}\hat f$, for $f\in L_{1}(-\infty,\infty)\cap L_{2}(-\infty,\infty)$ (cf.\ \eqref{eq:I-1.2}).

\section{Plancherel's theorem}
\label{sec:II-2}

\subsection*{Fourier transforms on $\mathcal S$}
\noindent If $\mathcal S$ denotes Schwartz's space (as in \S\ref{sec:I-3}, Ch.~I), and $f\in\mathcal S$, we define the Fourier transform $F[f]$ of $f$ by the relation
\begin{equation}\tag{2.1}\label{eq:II-2.1}
\begin{aligned}
F[f] &= \frac{1}{\sqrt{2\pi}}\,\hat f,\qquad\text{or}\qquad
F[f](\alpha)\\
&= \frac{1}{\sqrt{2\pi}}\,\hat f(\alpha)\\
&= \frac{1}{\sqrt{2\pi}}\int_{-\infty}^{\infty}f(x)e^{i\alpha x}\,dx,\quad-\infty<\alpha<\infty.
\end{aligned}
\end{equation}
The map $f\mapsto F[f]$ is a linear map of $\mathcal S$ onto itself, from what we have already seen (in \S\ref{sec:I-8}, Ch.~I). If we define $\check F[f]$ by the relation
\begin{equation}\tag{2.2}\label{eq:II-2.2}
\begin{aligned}
&\check F[f](\alpha) = F[f](-\alpha) = \frac{1}{\sqrt{2\pi}}\,\hat f(-\alpha)\\
&= \frac{1}{\sqrt{2\pi}}  \int_{-\infty}^{\infty}f(x)e^{-i\alpha x}\,dx,\quad - \infty < \alpha < \infty,
\end{aligned}
\end{equation}
then (by \eqref{eq:I-8.12}, Ch.~I, p.~50) we have
\begin{equation}\tag{2.3}\label{eq:II-2.3}
\check F\bigl[\check F[f]\bigr]=f.
\end{equation}
Further if $f,g\in\mathcal S$, then (by \eqref{eq:I-8.13}, Ch.~I, p.~50) we have (cf.\ \eqref{eq:I-1.1})
\begin{equation}\tag{2.4}\label{eq:II-2.4}
(f,g)=\bigl(F[f],F[g]\bigr);\qquad \|f\|_{2}=\|F[f]\|_{2};
\end{equation}
and
\begin{equation}\tag{2.5}\label{eq:II-2.5}
F[f*g]=\sqrt{2\pi}\,F[f]\cdot F[g],
\end{equation}
where the star ``*'' denotes convolution (see \eqref{eq:I-8.14}, Ch.~I, p.~50).

\subsection*{Fourier transforms on $L_{2}(-\infty,\infty)$}
\noindent Since $\mathcal S$ is a dense subset of $L_{2}(-\infty,\infty)$, if $f\in L_{2}(-\infty,\infty)$, there exists a sequence $(f_{n})$ of functions belonging to $\mathcal S$, such that $\|f-f_{n}\|_{2}\to 0$, as $n\to\infty$. Hence $\|f_{m}-f_{n}\|_{2}\to 0$, as $m,n\to\infty$, but by \eqref{eq:II-2.4}, and the linearity of the map $f_{n}\mapsto F[f_{n}]$, $\|F[f_{m}]-F[f_{n}]\|_{2}\to 0$, as $m,n\to\infty$. Since the space $L_{2}(-\infty,\infty)$ is complete, there exists a function $F\in L_{2}(-\infty,\infty)$, such that $\|F[f_{n}]-F\|_{2}\to 0$, as $n\to\infty$. The function $F$ is defined almost everywhere on $(-\infty,\infty)$. We call $F$ the Fourier transform of $f\in L_{2}(-\infty,\infty)$, and denote it by $F[f]$.

We note that $F$ does not depend on the approximating sequence $(f_{n})$. Let $(g_{n})$ be another sequence of functions belonging to $\mathcal S$, such that $\|g_{n}-f\|_{2}\to 0$. Then $f_{n}-g_{n}=f_{n}-f+f-g_{n}$, and $\|f_{n}-g_{n}\|_{2}\le \|f_{n}-f\|_{2}+\|g_{n}-f\|_{2}\to 0$, as $n\to\infty$. But, by \eqref{eq:II-2.4}, $\|F[f_{n}]-F[g_{n}]\|_{2}=\|f_{n}-g_{n}\|_{2}\to 0$, hence $\|F[g_{n}]-F\|_{2}\le \|F-F[f_{n}]\|_{2}+\|F[f_{n}]-F[g_{n}]\|_{2}\to 0$, as $n\to\infty$.

We are thus led to the following

\begin{nstmt}{2.6}\label{eq:II-2.6}
\textbf{Definition}. If $f\in L_{2}(-\infty,\infty)$, the Fourier transform $F[f]$ of $f$ is defined to be the limit, in the $L_{2}$-norm, of the sequence $\{F[f_{n}]\}$ of Fourier transforms, of any sequence $(f_{n})$ of functions belonging to $\mathcal S$, such that $f_{n}$ converges in the $L_{2}$-norm to the given function $f\in L_{2}(-\infty,\infty)$, as $n\to\infty$. The function $F[f]$ is defined almost everywhere on $(-\infty,\infty)$, and belongs to $L_{2}(-\infty,\infty)$.
\end{nstmt}

The definition makes it clear that the map $f\mapsto F[f]$ of $L_{2}(-\infty,\infty)$ into $L_{2}(-\infty,\infty)$ is linear. That is to say, if $f_{1},f_{2}\in L_{2}(-\infty,\infty)$, and $a,b\in\C$, then
\begin{equation}\tag{2.7}\label{eq:II-2.7}
F[a f_{1}+b f_{2}]=a\,F[f_{1}]+b\,F[f_{2}].
\end{equation}
We note further that for any $f\in L_{2}(-\infty,\infty)$, we have
\begin{equation}\tag{2.8}\label{eq:II-2.8}
\|f\|_{2}=\|F[f]\|_{2},
\end{equation}
since, by definition \eqref{eq:II-2.6}, and \eqref{eq:II-2.4}, we have
\begin{equation}\tag{2.9}\label{eq:II-2.9}
\|f\|_{2}=\lim_{n\to\infty}\|f_{n}\|_{2}
=\lim_{n\to\infty}\|F[f_{n}]\|_{2}=\|F[f]\|_{2}.
\end{equation}
If $f_{1},f_{2}\in L_{2}(-\infty,\infty)$, then we can apply \eqref{eq:II-2.8} to the function $f=f_{1}+f_{2}$, and noting that $\|f_{1}\|_{2}=\|F[f_{1}]\|_{2}$, $\|f_{2}\|_{2}=\|F[f_{2}]\|_{2}$, deduce that
\begin{equation}\tag{2.10}\label{eq:II-2.10}
(f_{1},f_{2})=\bigl(F[f_{1}],F[f_{2}]\bigr),\qquad\text{for }f_{1},f_{2}\in L_{2}(-\infty,\infty).
\end{equation}
We can similarly define $\check F[f]$ for any $f\in L_{2}(-\infty,\infty)$ by extending definition \eqref{eq:II-2.2} from $\mathcal S$ to $L_{2}(-\infty,\infty)$. Thus given any $f\in L_{2}(-\infty,\infty)$, there exists a sequence $(f_{n})$, $f_{n}\in\mathcal S$ for $n=1,2,\dots$, such that $\|f-f_{n}\|_{2}\to 0$, as $n\to\infty$. The sequence $(\check F[f_{n}])$, where $\check F[f_{n}]\in\mathcal S$, converges in the $L_{2}$-norm to a function in $L_{2}(-\infty,\infty)$, which we define to be $\check F[f]$. It is defined almost everywhere on $(-\infty,\infty)$, and we have
\begin{equation}\tag{2.11}\label{eq:II-2.11}
\check F[f](\alpha)=F[f](-\alpha)
\end{equation}
for almost all $\alpha\in(-\infty,\infty)$.

The ``inversion formula'' holds almost everywhere, without any special artifice, namely
\begin{equation*}\tag{2.11}\label{eq:II-2.11b}
\check F\bigl[F[f]\bigr]=f=F\bigl[\check F[f]\bigr],\qquad\text{for }f\in L_{2}(-\infty,\infty),
\end{equation*}
so that the map $f\mapsto F[f]$ is a linear map of $L_{2}(-\infty,\infty)$ onto itself; it is also isometric.

We shall now show that $F[f]$ thus defined for $f\in L_{2}(-\infty,\infty)$ is related to the integral defining $\frac{1}{\sqrt{2\pi}}\hat f$ in the case of $f\in L_{1}(-\infty,\infty)$.

Let $f\in L_{2}(-\infty,\infty)$, and let $f(x)=0$ for $|x|\ge A>0$. Then we shall see that
\begin{equation}\tag{2.12}\label{eq:II-2.12}
\begin{aligned}
&F[f](\alpha) = \frac{\hat f(\alpha)}{\sqrt{2\pi}}\\
&= \frac{1}{\sqrt{2\pi}}  \int_{-\infty}^{\infty}f(x)e^{i\alpha x}\,dx,  \qquad\text{for almost all }\alpha\in(-\infty,\infty).
\end{aligned}
\end{equation}
If we denote the integral on the right-hand side of \eqref{eq:II-2.12} by $f^{*}(\alpha)$, we shall show that $F[f](\alpha)=f^{*}(\alpha)$ for almost all $\alpha\in(-\infty,\infty)$.

There exists a sequence $(f_{n})$, $n=1,2,\dots$, with $f_{n}\in\mathcal S$, such that every member of the sequence vanishes outside the interval $(-A,A)$, and $\|f-f_{n}\|_{2}\to 0$, as $n\to\infty$. That implies, by \eqref{eq:II-2.7} and \eqref{eq:II-2.8}, that $\|F[f]-F[f_{n}]\|_{2}\to 0$, and, in particular,
\begin{equation}\tag{2.13}\label{eq:II-2.13}
\int_{-R}^{R}\bigl|F[f_{n}](x)-F[f](x)\bigr|^{2}\,dx\to 0,\qquad\text{as }n\to\infty,
\end{equation}
for each $R>0$. But
\[
\begin{aligned}
&\bigl|F[f_{n}](\alpha) - f^{*}(\alpha)\bigr|\\
&\le \int_{-A}^{A}  \bigl|f_{n}(x) - f(x)\bigr|\,dx \le \Bigl(2A\int_{-A}^{A}\bigl|f_{n}(x)-f(x)\bigr|^{2}\,dx\Bigr)^{1/2},
\end{aligned}
\]
hence $F[f_{n}]\to f^{*}$ uniformly over every finite interval, and, in particular,
\begin{equation}\tag{2.14}\label{eq:II-2.14}
\int_{-R}^{R}\bigl|F[f_{n}](x)-f^{*}(x)\bigr|^{2}\,dx\to 0,\qquad\text{as }n\to\infty.
\end{equation}
A comparison of \eqref{eq:II-2.13} and \eqref{eq:II-2.14} shows that $f^{*}=F[f]$ almost everywhere on $(-R,R)$, for each $R>0$, and hence almost everywhere on $(-\infty,\infty)$, thus proving \eqref{eq:II-2.12}.

Finally let $f\in L_{2}(-\infty,\infty)$ without any further condition. We define the function $f_{R}$, for each $R>0$, by the requirement
\[
f_{R}(x)=
\begin{cases}
f(x), & |x|<R,\\
0, & |x|\ge R.
\end{cases}
\]
Then $f_{R}\in L_{1}(-\infty,\infty)\cap L_{2}(-\infty,\infty)$, and by what has just been proved in \eqref{eq:II-2.12}, we have
\begin{equation}\tag{2.15}\label{eq:II-2.15}
F[f_{R}](\alpha)=\frac{1}{\sqrt{2\pi}}\int_{-R}^{R}f(x)e^{i\alpha x}\,dx,\qquad-\infty<\alpha<\infty.
\end{equation}
On the other hand, since $f\mapsto F[f]$ is a linear map of $L_{2}(-\infty,\infty)$ onto itself, we have, by \eqref{eq:II-2.8},
\begin{equation}\tag{2.16}\label{eq:II-2.16}
\bigl\|F[f]-F[f_{R}]\bigr\|_{2}
=\bigl\|F[f-f_{R}]\bigr\|_{2}
=\|f-f_{R}\|_{2}\to 0,\qquad\text{as }R\to\infty.
\end{equation}
Hence for each $f\in L_{2}(-\infty,\infty)$, the integral $F[f_{R}]$ given by \eqref{eq:II-2.15} converges in the $L_{2}$-norm to $F[f]\in L_{2}(-\infty,\infty)$, as $R\to\infty$. This implies, by Weyl's formulation of the Riesz--Fischer theorem, that there exists a sequence $(R_{k})$, with $R_{k}>0$ for $k=1,2,\dots$, such that
\[
\frac{1}{\sqrt{2\pi}}\int_{-R_{k}}^{R_{k}}f(x)e^{ix\alpha}\,dx\to F[f](\alpha),
\]
as $R_{k}\to\infty$, for almost every $\alpha\in(-\infty,\infty)$. In particular, if the integral
\begin{equation}\tag{2.17}\label{eq:II-2.17}
\frac{1}{\sqrt{2\pi}}\int_{-\infty}^{\infty}f(x)e^{ix\alpha}\,dx
\end{equation}
exists as a Cauchy principal value for almost all $\alpha\in(-\infty,\infty)$, it equals $F[f]$. And if $f\in L_{1}(-\infty,\infty)\cap L_{2}(-\infty,\infty)$, then
\begin{equation}\tag{2.18}\label{eq:II-2.18}
F[f](\alpha)=\frac{1}{\sqrt{2\pi}}\int_{-\infty}^{\infty}f(x)e^{i\alpha x}\,dx,\qquad-\infty<\alpha<\infty.
\end{equation}
We subsume the results of \eqref{eq:II-2.6}--\eqref{eq:II-2.11} (and \eqref{eq:II-2.11b}), \eqref{eq:II-2.17} and \eqref{eq:II-2.18} under the following

\begin{theorem}[Plancherel]\label{thm:II-1}
If $f\in L_{2}(-\infty,\infty)$, then there exists a function $F[f]\in L_{2}(-\infty,\infty)$, designated the Fourier transform of $f$, such that, for any real $\alpha$,
\begin{equation}\tag{2.19}\label{eq:II-2.19}
\frac{1}{\sqrt{2\pi}}\int_{-R}^{R}f(x)e^{i\alpha x}\,dx\to F[f](\alpha),\qquad\text{in the $L_{2}$-norm, as $R\to\infty$},
\end{equation}
and
\begin{equation}\tag{2.20}\label{eq:II-2.20}
\frac{1}{\sqrt{2\pi}}\int_{-R}^{R}F[f](\alpha)e^{-i\alpha x}\,d\alpha\to f(x),\qquad\text{in the $L_{2}$-norm, as $R\to\infty$},
\end{equation}
with
\begin{equation}\tag{2.21}\label{eq:II-2.21}
\|f\|_{2}=\|F[f]\|_{2}.
\end{equation}
Every function $f\in L_{2}(-\infty,\infty)$ is the Fourier transform of a unique element of $L_{2}(-\infty,\infty)$.
\end{theorem}

As in the $L_{1}$-case (see Ch.~I, \eqref{eq:I-1.8}, \eqref{eq:I-1.9}), we have for $f\in L_{2}(-\infty,\infty)$, and any real number $a$, the relations
\begin{align}\tag{2.22}\label{eq:II-2.22}
F[f(\cdot+a)](y)&=e^{-iya}\,F[f](y), & F[\overline{f}(\cdot)](y)&=\overline{F[f]}(-y);
\end{align}
while
\begin{align}\tag{2.23}\label{eq:II-2.23}
\check F[f(\cdot+a)](y)&=e^{iya}\,\check F[f](-y), & \check F[\overline{f}(\cdot)](y)&=\overline{\check F[f]}(y).
\end{align}
To indicate the reasoning involved, let us consider the first relation in \eqref{eq:II-2.22}. It follows from the fact that
\[
\begin{aligned}
&\int_{-R}^{R}f(x+a)e^{iyx}\,dx\\
&= \int_{-R+a}^{R+a}f(x)e^{iy(x-a)}\,dx = e^{-iya}  \int_{-R+a}^{R+a}f(x)e^{iyx}\,dx,
\end{aligned}
\]
where
\[
\int_{-R+a}^{-R}f(x)e^{iyx}\,dx\to 0,\qquad\text{and}\qquad
\int_{R}^{R+a}f(x)e^{iyx}\,dx\to 0,
\]
in the $L_{2}$-norm, as $R\to\infty$, because of \eqref{eq:II-2.9}.

For $f,g\in L_{2}(-\infty,\infty)$, we have already shown in \eqref{eq:II-2.10} that \eqref{eq:II-2.8} implies that
\begin{equation}\tag{2.24}\label{eq:II-2.24}
\int_{-\infty}^{\infty}f(x)\,\overline{g(x)}\,dx
=\int_{-\infty}^{\infty}F[f](x)\,\overline{F[g](x)}\,dx.
\end{equation}
On using \eqref{eq:II-2.22} and \eqref{eq:II-2.23} in \eqref{eq:II-2.24}, we get the following relations:
\begin{equation}\tag{2.25}\label{eq:II-2.25}
\int_{-\infty}^{\infty}f(x)\,g(-x)\,dx
=\int_{-\infty}^{\infty}F[f](y)\,F[g](y)\,dy;
\end{equation}
\begin{equation}\tag{2.26}\label{eq:II-2.26}
\int_{-\infty}^{\infty}f(x)\,g(a-x)\,dx
=\int_{-\infty}^{\infty}F[f](x)\,F[g](x)\,e^{-iax}\,dx;
\end{equation}
and
\begin{equation}\tag{2.27}\label{eq:II-2.27}
\int_{-\infty}^{\infty}f(t)\,g(t)\,e^{ixt}\,dt
=\int_{-\infty}^{\infty}F[f](t)\,F[g](x-t)\,dt,
\end{equation}
all the integrals being absolutely convergent.

We may look upon \eqref{eq:II-2.26} as the $L_{2}$-analogue of formula \eqref{eq:I-1.2} (or \eqref{eq:I-2.2}) of Chapter~I.

If we choose
\[
g(x)=g_{a,b}(x)=
\begin{cases}
1, & a<x<b,\\
0, & x\le a,\quad x\ge b,
\end{cases}
\]
then $g\in L_{1}(-\infty,\infty)\cap L_{2}(-\infty,\infty)$, and
\[
F[g](y)=\frac{1}{\sqrt{2\pi}}\int_{a}^{b}e^{ixy}\,dx
=\frac{1}{\sqrt{2\pi}}\frac{e^{iby}-e^{iay}}{iy}\in L_{2}(-\infty<y<\infty),
\]
and \eqref{eq:II-2.27}, with $x=0$, gives
\begin{equation}\tag{2.28}\label{eq:II-2.28}
\int_{a}^{b}f(t)\,dt
=\frac{1}{\sqrt{2\pi}}\int_{-\infty}^{\infty}F[f](y)\,\frac{e^{-iby}-e^{-iay}}{-iy}\,dy,\qquad a<b.
\end{equation}
On the other hand, we may take
\[
F[g](y)=F[g_{a,b}](y)=
\begin{cases}
1, & a<y<b,\\
0, & y\le a,\quad y\ge b.
\end{cases}
\]
Then
\[
\begin{aligned}
g(x) &= \frac{1}{\sqrt{2\pi}}\int_{-\infty}^{\infty}F[g_{a,b}](y)e^{-ixy}\,dy\\
&= \frac{1}{\sqrt{2\pi}}\int_{a}^{b}e^{-ixy}\,dy\\
&= \frac{1}{\sqrt{2\pi}}\frac{e^{-ibx}-e^{-iax}}{-ix}\in L_{2}(-\infty<x<\infty),
\end{aligned}
\]
and \eqref{eq:II-2.25} gives
\begin{equation}\tag{2.29}\label{eq:II-2.29}
\int_{a}^{b}F[f](x)\,dx
=\frac{1}{\sqrt{2\pi}}\int_{-\infty}^{\infty}f(x)\,\frac{e^{ibx}-e^{iax}}{ix}\,dx,\qquad a<b.
\end{equation}

\begin{remark}
The operator $F$ has been defined on $L_{1}(-\infty,\infty)$ in Chapter~I, \eqref{eq:I-1.2}, and on $L_{2}(-\infty,\infty)$ in Plancherel's theorem. The definition can be extended to the space $L_{1}(-\infty,\infty)+L_{2}(-\infty,\infty)$ consisting of all functions $f$ of the form $f=f_{1}+f_{2}$, where $f_{1}\in L_{1}(-\infty,\infty)$, $f_{2}\in L_{2}(-\infty,\infty)$.

For such an $f$ we define $F[f]=F[f_{1}]+F[f_{2}]$. This definition does not depend on the particular decomposition $f=f_{1}+f_{2}$. For if $f=g_{1}+g_{2}$, $g_{1}\in L_{1}(-\infty,\infty)$, $g_{2}\in L_{2}(-\infty,\infty)$ is a different decomposition, then $f_{1}-g_{1}=g_{2}-f_{2}\in L_{1}(-\infty,\infty)\cap L_{2}(-\infty,\infty)$ on which, by \eqref{eq:II-2.18}, the $L_{1}$-definition and the $L_{2}$-definition of $F$ coincide, so that $F[f_{1}-g_{1}]=F[f_{1}]-F[g_{1}]=F[g_{2}]-F[f_{2}]$, and we have $F[f_{1}]+F[f_{2}]=F[g_{1}]+F[g_{2}]$. Thus the operator $F$ is defined on the space $L_{1}(-\infty,\infty)+L_{2}(-\infty,\infty)$, which, in fact, contains all the spaces $L_{p}(-\infty,\infty)$, for $1\le p\le 2$. For if $g\in L_{p}(-\infty,\infty)$, $1<p<2$, then $g=g_{1}+g_{2}$, where $g_{1}\in L_{1}(-\infty,\infty)$, and $g_{2}\in L_{2}(-\infty,\infty)$. We have only to define
\[
g_{1}(x)=
\begin{cases}
g(x), & |g(x)|\ge 1,\\
0, & |g(x)|<1,
\end{cases}
\qquad\text{and}\qquad
g_{2}(x)=
\begin{cases}
g(x), & |g(x)|<1,\\
0, & |g(x)|\ge 1,
\end{cases}
\]
so that
\[
\begin{aligned}
&\int_{-\infty}^{\infty}|g_{1}(x)|\,dx\\
&= \int_{|g(x)|\ge 1}|g(x)|\,dx \le \int_{|g(x)|\ge 1}|g(x)|^{p}\,dx \le \int_{-\infty}^{\infty}|g(x)|^{p}\,dx < \infty,
\end{aligned}
\]
while $g_{2}$ is bounded, and belongs to $L_{p}(-\infty,\infty)$, since
\[
\int_{-\infty}^{\infty}|g_{2}(x)|^{p}\,dx
=\int_{|g(x)|<1}|g(x)|^{p}\,dx<\infty,
\]
and hence
\[
\int_{-\infty}^{\infty}|g_{2}(x)|^{2}\,dx
=\int_{|g(x)|<1}|g(x)|^{p+(2-p)}\,dx
=\int_{|g(x)|<1}|g(x)|^{p}\,dx<\infty.
\]
The operator $F$ so defined on $L_{p}(-\infty,\infty)$, $1<p<2$, may be called the Fourier transform as well.
\end{remark}

\subsection*{Convergence and summability}
\noindent The problem of convergence of the ``Fourier integral''
\begin{equation}\tag{3.1}\label{eq:II-3.1}
S_{R}(f;x)\equiv \frac{1}{\sqrt{2\pi}}\int_{-R}^{R}F[f](\alpha)e^{-i\alpha x}\,d\alpha,\qquad R>0,\quad f\in L_{2}(-\infty,\infty),
\end{equation}
$-\infty<x<\infty$, as $R\to\infty$, can be handled in the same way as in the $L_{1}$-case, (Ch.~I, \S\S\ref{sec:I-4}, \ref{sec:I-7}), once the basic formula (cf.\ \eqref{eq:I-4.2}, Ch.~I)
\begin{equation}\tag{3.2}\label{eq:II-3.2}
S_{R}(f;x)-f(x)=\frac{2}{\pi}\int_{0}^{\infty}g_{x}(t)\,\frac{\sin Rt}{t}\,dt,
\end{equation}
where $g_{x}(t)=\frac{1}{2}[f(x+t)+f(x-t)-2f(x)]$, is established, which can be done by using Plancherel's theorem.

For let
\[
F[\varphi](y)=
\begin{cases}
e^{-iyx}, & |y|<R,\\
0, & |y|\ge R,
\end{cases}
\qquad R>0,\quad -\infty<x<\infty.
\]
Then we have
\[
\begin{aligned}
\varphi(u) &= \frac{1}{\sqrt{2\pi}}\int_{-\infty}^{\infty}F[\varphi](y)e^{-iyu}\,dy\\
&= \frac{1}{\sqrt{2\pi}}\int_{-R}^{R}e^{-iy(x+u)}\,dy\\
&= \Bigl(\frac{2}{\pi}\Bigr)^{1/2}\frac{\sin R(x+u)}{x+u},
\end{aligned}
\]
and formula \eqref{eq:II-2.25} gives
\[
\begin{aligned}
&\int_{-R}^{R}F[f](\alpha)e^{-i\alpha x}\,d\alpha = \int_{-\infty}^{\infty}f(u)\varphi(-u)\,du\\
&= \Bigl(\frac{2}{\pi}\Bigr)^{1/2}\int_{-\infty}^{\infty}f(x-t)\frac{\sin Rt}{t}\,dt\\
&= \Bigl(\frac{2}{\pi}\Bigr)^{1/2}\int_{0}^{\infty}[f(x+t)+f(x-t)]\frac{\sin Rt}{t}\,dt,
\end{aligned}
\]
which leads to \eqref{eq:II-3.2}. The following analogue of Theorem~5 of Chapter~I is a consequence, the proof being similar.

\begin{theorem}\label{thm:II-2}
If $f\in L_{2}(-\infty,\infty)$, and $f$ is of bounded variation in a neighbourhood of the point $x\in(-\infty,\infty)$, then
\[
\begin{aligned}
&\lim_{R\to\infty}S_{R}(f;x)\\
&= \lim_{R\to\infty}\frac{1}{\sqrt{2\pi}}  \int_{-R}^{R}F[f](\alpha)e^{-i\alpha x}\,d\alpha = \frac{1}{2}  \bigl[f(x+0)+f(x-0)\bigr].
\end{aligned}
\]
\end{theorem}

The problem of $(C,1)$ summability, for example, can again be handled without any new difficulty once the basic formula (cf.\ \eqref{eq:I-7.2}, Ch.~I)
\begin{equation}\tag{3.3}\label{eq:II-3.3}
\begin{aligned}
&\frac{1}{\sqrt{2\pi}}  \int_{-R}^{R}F[f](\alpha)  \Bigl(1-\frac{|\alpha|}{R}\Bigr)e^{-i\alpha x}\,d\alpha - f(x)\\
&= \frac{1}{\pi}  \int_{0}^{\infty}g_{x}(t)\,\frac{\sin^{2}(Rt/2)}{R(t/2)^{2}}\,dt
\end{aligned}
\end{equation}
is established. Let
\[
F[\varphi](y)=
\begin{cases}
\bigl(1-\frac{|y|}{R}\bigr)e^{-ixy}, & |y|\le R,\\
0, & |y|>R.
\end{cases}
\]
Then we have (cf.\ Example~2, \S\ref{sec:I-1}, Ch.~I)
\[
\varphi(u)=\frac{1}{\sqrt{2\pi}}\int_{-R}^{R}\Bigl(1-\frac{|y|}{R}\Bigr)e^{-i(x+u)y}\,dy
=\frac{4}{\sqrt{2\pi}}\frac{\sin^{2}[R(x+u)/2]}{R(x+u)^{2}},
\]
so that formula \eqref{eq:II-2.25} again gives
\[
\begin{aligned}
&\frac{1}{\sqrt{2\pi}}\int_{-R}^{R}F[f](\alpha)\Bigl(1-\frac{|\alpha|}{R}\Bigr)e^{-i\alpha x}\,d\alpha\\
&= \frac{2}{\pi}\int_{-\infty}^{\infty}f(x-t)\frac{\sin^{2}(Rt/2)}{Rt^{2}}\,dt\\
&= \frac{1}{2\pi}\int_{0}^{\infty}[f(x+t)+f(x-t)]\frac{\sin^{2}(Rt/2)}{R(t/2)^{2}}\,dt,
\end{aligned}
\]
which leads to \eqref{eq:II-3.3}. Following the same lines of proof as in Theorem~10 of Chapter~I, we obtain, for example,

\begin{theorem}\label{thm:II-3}
If $f\in L_{2}(-\infty,\infty)$, then
\[
\lim_{R\to\infty}\frac{1}{\sqrt{2\pi}}\int_{-R}^{R}F[f](\alpha)\Bigl(1-\frac{|\alpha|}{R}\Bigr)e^{-i\alpha x}\,d\alpha=f(x),
\]
for almost all $x\in(-\infty,\infty)$.
\end{theorem}

We note that formula \eqref{eq:II-2.24} plays the same r\^ole in the $L_{2}$-theory as the composition rule (cf.\ \eqref{eq:I-1.13}, Ch.~I) did in the $L_{1}$-case. The problem of inversion in $L_{2}(-\infty,\infty)$ does not arise in the same form as in $L_{1}(-\infty,\infty)$, and the theorems on point-wise summability of the Fourier integral, while true, do not carry the same import.

To consider the question of summability in the $L_{2}$-norm, we need an extension of Theorem~2 of Chapter~I, on the convolution of two functions in $L_{1}(-\infty,\infty)$.

\begin{nstmt}{3.4}\label{eq:II-3.4}
If $f\in L_{1}(-\infty,\infty)$, and $g\in L_{p}(-\infty,\infty)$, where $1\le p\le 2$, then for almost every $x\in(-\infty,\infty)$, the function $f(x-y)\,g(y)$ belongs to $L_{1}(-\infty<y<\infty)$, and if we define $f*g$ by the relation
\end{nstmt}
\[
(f*g)(x)=\int_{-\infty}^{\infty}f(x-y)g(y)\,dy,
\]
then
\begin{equation}\tag{3.5}\label{eq:II-3.5}
\|f*g\|_{p}\le \|f\|_{1}\,\|g\|_{p},
\end{equation}
so that $f*g\in L_{p}(-\infty,\infty)$.

This is proved in the same way as Lemma~2.1 and Theorem~2, of Chapter~I, except that one has to use H\"older's inequality instead of Schwarz's.

In particular, if $p=2$, and $h=f*g$, then $h\in L_{2}(-\infty,\infty)$, so that $F[h]$ is defined, as well as $F[f]$ (cf.\ \eqref{eq:I-1.2}, Ch.~I), and $F[g]$. We shall now prove

\begin{theorem}\label{thm:II-4}
If $f\in L_{1}(-\infty,\infty)$, and $g\in L_{2}(-\infty,\infty)$, and $h=f*g$, then $h\in L_{2}(-\infty,\infty)$, and
\begin{equation*}\tag{3.6}\label{eq:II-3.6}
F[h](x)=\sqrt{2\pi}\,F[f](x)\,F[g](x),
\end{equation*}
for almost all $x\in(-\infty,\infty)$.
\end{theorem}
\begin{proof}
Obviously $F[f]$ is bounded; in fact, $\sqrt{2\pi}\,|F[f](\alpha)|\le \|f\|_{1}=M<\infty$, and if $(g_{n})$ is a sequence of functions belonging to Schwartz's space $\mathcal S$ (cf.\ \S\ref{sec:I-3}, Ch.~I), such that $\|g-g_{n}\|_{2}\to 0$, as $n\to\infty$, and if we set $h_{n}=f*g_{n}$, then $h_{n}\in L_{1}(-\infty,\infty)\cap L_{2}(-\infty,\infty)$, and $F[h_{n}]=\sqrt{2\pi}\,F[f]\,F[g_{n}]$ by Theorem~2 of Chapter~I. Further, because of \eqref{eq:II-3.5}, $\|h-h_{n}\|_{2}=\|f*(g-g_{n})\|_{2}\to 0$, as $n\to\infty$. Since $\|F[h]-F[h_{n}]\|_{2}=\|h-h_{n}\|_{2}\to 0$, as $n\to\infty$, and
\[
\begin{aligned}
& \|F[h_{n}] - \sqrt{2\pi}\,F[f]\,F[g]\|_{2} = \|F[f]\,F[g_{n}] - F[f]\,F[g]\|_{2}\\
&\le M\,\|F[g_{n}] - F[g]\|_{2}\\
&= M\,\|g-g_{n}\|_{2}\to 0,
\end{aligned}
\]
as $n\to\infty$, we obtain \eqref{eq:II-3.6}.
\end{proof}

The question of summability in the $L_{2}$-norm is answered by the following

\begin{theorem}\label{thm:II-5}
Let $f\in L_{2}(-\infty,\infty)$, $K\in L_{2}(-\infty,\infty)$, and $F[K]=H\in L_{1}(-\infty,\infty)\cap L_{2}(-\infty,\infty)$, $H$ even. For $R>0$, let $H_{R}(y)=R\,H(Ry)$. Then we have
\begin{equation*}\tag{3.7}\label{eq:II-3.7}
\Bigl\|\frac{1}{\sqrt{2\pi}}\,(f*H_{R})-f\Bigr\|_{2}\to 0,\qquad\text{as }R\to\infty.
\end{equation*}
\end{theorem}
\begin{proof}
By formula \eqref{eq:II-2.25}, we have
\[
\begin{aligned}
&\int_{-\infty}^{\infty}F[f](y)\,K\Bigl(\frac{y}{R}\Bigr)e^{-iyx}\,dy = \int_{-\infty}^{\infty}f(x+y)\,H_{R}(y)\,dy\\
&= \int_{-\infty}^{\infty}f(x-y)\,H_{R}(y)\,dy\\
&= (f*H_{R})(x).
\end{aligned}
\]
By \eqref{eq:II-3.5} we see that $f*H_{R}\in L_{2}(-\infty,\infty)$, for each $R>0$. By using the properties of the $L_{2}$-modulus of continuity (see \eqref{eq:I-1.17}, Ch.~I), we deduce as before (cf.\ Th.~16, Ch.~I) that
\[
\Bigl\|\frac{1}{\sqrt{2\pi}}\,(f*H_{R})-f\Bigr\|_{2}\to 0,\qquad\text{as }R\to\infty.
\]
\end{proof}
\begin{remark}
Clearly we can take for $K(\alpha)$ the Abel kernel $e^{-|\alpha|}$, and the Gauss kernel $e^{-\alpha^{2}}$, as we did in Chapter~I.
\end{remark}

\setcounter{section}{3}
\section{The closure of translations}
\label{sec:II-4}

Let $f\in L_{2}(-\infty,\infty)$, and let $\overline{S}_{f}$ denote the closure, in the $L_{2}$-norm, of the set $S_{f}$ of all ``translations'' of the form $\sum_{k=1}^{m}c_{k}f(\cdot+t_{k})$, where $t_{k}$ is real, and $c_{k}$ complex (cf.\ \S\ref{sec:I-12}, Ch.~I). Then we have the following

\begin{theorem}[Wiener]\label{thm:II-6}
If $f\in L_{2}(-\infty,\infty)$, with $F[f](\alpha)\neq 0$ for almost all $\alpha\in(-\infty,\infty)$, then $\overline{S}_{f}=L_{2}(-\infty,\infty)$.
\end{theorem}

This can be obtained as a special case of another theorem which provides a sufficient condition for a function in $L_{2}(-\infty,\infty)$ to belong to $\overline{S}_{f}$.

Let $\mathbb R_{1}$ denote the set of all real numbers, and let $\varphi=F[f]$ where $f\in L_{2}(-\infty,\infty)$. Let $E_{\varphi}$ denote a set in $\mathbb R_{1}$ with the property that, except for null sets (i.e.\ measurable sets of Lebesgue measure zero), $\varphi(\alpha)\neq 0$ for $\alpha\in E_{\varphi}$, and $\varphi(\alpha)=0$ for $\alpha\in\mathbb R_{1}-E_{\varphi}$. The set $E_{\varphi}$ is, of course, measurable.

If $F(x)=\sum_{k=1}^{m}c_{k}\,f(x+t_{k})$, $F$ is not identically zero, and $\Phi=F[F]$, then
\begin{equation}\tag{4.1}\label{eq:II-4.1}
\Phi(\alpha)=\{c_{1}e^{-i\alpha t_{1}}+\dots+c_{m}e^{-i\alpha t_{m}}\}\,\varphi(\alpha),
\end{equation}
and clearly $E_{\Phi}=E_{\varphi}$, except for null sets, and if $g\in\overline{S}_{f}$, with $F[g]=\psi$, then
\begin{equation}\tag{4.2}\label{eq:II-4.2}
E_{\psi}<E_{\varphi},
\end{equation}
which is defined to mean that $E_{\psi}=E_{1}\cup E_{2}$, where $E_{1}\subset E_{\psi}$ and $E_{2}$ is a null set. But the converse is also true, as shown by the following

\begin{theorem}\label{thm:II-7}
If $f,g\in L_{2}(-\infty,\infty)$, and $F[f]=\varphi$, $F[g]=\psi$, and if
\begin{equation*}\tag{4.2}\label{eq:II-4.2b}
E_{\psi}<E_{\varphi},
\end{equation*}
then $g\in\overline{S}_{f}$.
\end{theorem}
\begin{proof}
Clearly $\overline{S}_{f}$ is a closed, linear subspace of $L_{2}(-\infty,\infty)$, which is a separable Hilbert space. Given any $g\in L_{2}(-\infty,\infty)$, there exists an element $g^{0}\in\overline{S}_{f}$, such that
\begin{equation}\tag{4.3}\label{eq:II-4.3}
g=g^{0}+f^{0},
\end{equation}
with
\begin{equation}\tag{4.4}\label{eq:II-4.4}
(f^{0},h)=\int_{-\infty}^{\infty}f^{0}(x)\,\overline{h(x)}\,dx=0,
\end{equation}
for every $h\in\overline{S}_{f}$. If $\varphi^{0}=F[f^{0}]$, we have again
\begin{equation}\tag{4.5}\label{eq:II-4.5}
E_{\varphi^{0}}<E_{\varphi},
\end{equation}
because of \eqref{eq:II-4.1} and \eqref{eq:II-4.2}. If we now choose $h(x)=f(x+t)$, for a fixed, real $t$, in \eqref{eq:II-4.4}, then we have, by \eqref{eq:II-2.24} and \eqref{eq:II-2.22},
\[
\int_{-\infty}^{\infty}\varphi^{0}(\alpha)\,\overline{\varphi(\alpha)}\,e^{i\alpha t}\,d\alpha=0,
\]
where $\varphi^{0}\cdot\varphi\in L_{1}(-\infty,\infty)$. Hence, by Theorem~7 of Chapter~I,
\begin{equation}\tag{4.6}\label{eq:II-4.6}
\varphi^{0}(\alpha)\,\overline{\varphi(\alpha)}=0,
\end{equation}
for almost all $\alpha\in(-\infty,\infty)$. But we have (except for sets of measure zero) $\varphi^{0}(\alpha)=0$ for $\alpha\in\mathbb R_{1}-E_{\varphi}$, by \eqref{eq:II-4.5}, while $\varphi(\alpha)\neq 0$ for $\alpha\in E_{\varphi}$. Hence $\varphi^{0}(\alpha)=0$ for almost every $\alpha$. By Plancherel's theorem, $f^{0}(x)=0$, for almost every $x$, which gives, because of \eqref{eq:II-4.3}, $g(x)=g^{0}(x)$ for almost every $x$, and hence $g\in\overline{S}_{f}$.
\end{proof}

\begin{proof}[Proof of Theorem~\ref{thm:II-6}]
If $E_{\varphi}=\mathbb R_{1}$, then automatically we have $E_{\psi}<E_{\varphi}$, for every $g\in L_{2}(-\infty,\infty)$. Theorem~\ref{thm:II-7} then implies that $g\in\overline{S}_{f}$, for every $g\in L_{2}(-\infty,\infty)$.
\end{proof}

\setcounter{section}{4}
\section{Heisenberg's inequality}
\label{sec:II-5}

We shall now prove an $L_{2}$-analogue of what is referred to as Heisenberg's inequality, originally proved by Weyl under somewhat stronger assumptions.

\begin{theorem}\label{thm:II-8}
Let $f\in L_{2}(-\infty,\infty)$. Then we have
\begin{equation}\tag{5.1}\label{eq:II-5.1}
\int_{-\infty}^{\infty}x^{2}\,|f(x)|^{2}\,dx
\int_{-\infty}^{\infty}\alpha^{2}\,|F[f](\alpha)|^{2}\,d\alpha
\ge \frac{1}{4}\,\|f\|_{2}^{4},
\end{equation}
and the equality takes place only in case $f(x)=c\,e^{-kx^{2}}$, $k>0$, $c\in\R$.
\end{theorem}
\begin{proof}
We may assume that
\[
\int_{-\infty}^{\infty}x^{2}\,|f(x)|^{2}\,dx<\infty,\qquad
\int_{-\infty}^{\infty}\alpha^{2}\,|F[f](\alpha)|^{2}\,d\alpha<\infty,
\]
for otherwise \eqref{eq:II-5.1} is trivially true since neither term can be zero.

Let $f^{*}(\alpha)$ denote the inverse transform of $(-i\alpha\,F[f](\alpha))$, that is to say
\begin{equation}\tag{5.2}\label{eq:II-5.2}
f^{*}(\alpha)=\check F\bigl[-i\alpha\,F[f](\alpha)\bigr].
\end{equation}
By Plancherel's theorem, $f^{*}\in L_{2}(-\infty,\infty)$. The left-hand side of \eqref{eq:II-5.1} equals
\[
\begin{aligned}
& \int_{-\infty}^{\infty}x^{2}\,|f(x)|^{2}\,dx  \int_{-\infty}^{\infty}\alpha^{2}\,|F[f](\alpha)|^{2}\,d\alpha\\
&= \int_{-\infty}^{\infty}x^{2}\,|f(x)|^{2}\,dx  \int_{-\infty}^{\infty}  \bigl|(-i\alpha\,F[f](\alpha))\bigr|^{2}\,d\alpha\\
&= \int_{-\infty}^{\infty}x^{2}\,|f(x)|^{2}\,dx
\int_{-\infty}^{\infty}|f^{*}(\alpha)|^{2}\,d\alpha,
\end{aligned}
\]
by \eqref{eq:II-5.2} and Plancherel's theorem. By Schwarz's inequality, we have
\begin{equation}\tag{5.3}\label{eq:II-5.3}
\int_{-\infty}^{\infty}x^{2}\,|f(x)|^{2}\,dx
\int_{-\infty}^{\infty}|f^{*}(x)|^{2}\,dx
\ge \Bigl[\int_{-\infty}^{\infty}\frac{x\bigl(f^{*}\overline{f}+\overline{f^{*}}f\bigr)}{2}\,dx\Bigr]^{2},
\end{equation}
since $x f^{*}(x)\overline{f(x)}=\frac{1}{2}x\bigl(f^{*}\overline{f}+\overline{f^{*}}f\bigr)$. To prove the theorem it suffices therefore to show that
\begin{equation}\tag{5.4}\label{eq:II-5.4}
\Bigl[\int_{-\infty}^{\infty}x\bigl(f^{*}\overline{f}+\overline{f^{*}}f\bigr)\,dx\Bigr]^{2}=\|f\|_{2}^{4}.
\end{equation}

Let $f_{n}\in\mathcal S$ (Schwartz's space) be so chosen that
\begin{equation}\tag{5.5}\label{eq:II-5.5}
\lim_{n\to\infty}\int_{-\infty}^{\infty}(1+\alpha^{2})\,
\bigl|F[f_{n}](\alpha)-F[f](\alpha)\bigr|^{2}\,d\alpha=0.
\end{equation}
This is indeed possible, since $\mathcal S$ is a dense subset of $L_{2}(-\infty,\infty)$. Since $\alpha|F[f](\alpha)|\in L_{2}(-\infty,\infty)$ by assumption, and $F[f]\in L_{2}(-\infty,\infty)$, we have $(1+\alpha^{2})^{1/2}F[f](\alpha)\in L_{2}(-\infty,\infty)$. There exists a sequence $(g_{n})$, $g_{n}\in\mathcal S$, such that
\[
\int_{-\infty}^{\infty}\Bigl|g_{n}(\alpha)-\frac{F[f](\alpha)}{(1+\alpha^{2})^{1/2}}\Bigr|^{2}\,d\alpha\to 0,\qquad\text{as }n\to\infty,
\]
or
\[
\int_{-\infty}^{\infty}(1+\alpha^{2})\,
\Bigl|\frac{g_{n}(\alpha)}{(1+\alpha^{2})^{1/2}}-F[f](\alpha)\Bigr|^{2}\,d\alpha\to 0,\qquad\text{as }n\to\infty.
\]
Choose $f_{n}$ such that
\[
F[f_{n}](\alpha)=\frac{g_{n}(\alpha)}{(1+\alpha^{2})^{1/2}},
\]
noting that $F[f_{n}]\in\mathcal S$ implies that $f_{n}\in\mathcal S$ (cf.\ Ch.~I, \eqref{eq:I-3.7}).

Now by Plancherel's theorem, and \eqref{eq:II-5.5},
\[
\|f_{n}-f\|_{2}^{2}+\|f_{n}'-f^{*}\|_{2}^{2}
=\int_{-\infty}^{\infty}(1+\alpha^{2})\,
\bigl|F[f_{n}](\alpha)-F[f](\alpha)\bigr|^{2}\,d\alpha\to 0,
\]
as $n\to\infty$, where $f_{n}'$ denotes the derivative of $f_{n}$, so that $F[f_{n}'](\alpha)=-i\alpha\,F[f_{n}](\alpha)$, and $f^{*}$ is defined as in \eqref{eq:II-5.2}. It follows that
\begin{equation}\tag{5.6}\label{eq:II-5.6}
\|f_{n}-f\|_{2}\to 0,\qquad \|f_{n}'-f^{*}\|_{2}\to 0,\qquad\text{as }n\to\infty.
\end{equation}
Now
\[
\begin{aligned}
&\int_{-\infty}^{\infty}\bigl|F[f_{n}](\alpha)-\mathcal F[f](\alpha)\bigr|\,d\alpha
\le \Bigl[\int_{-\infty}^{\infty}\frac{d\alpha}{1+\alpha^{2}}\Bigr]^{1/2}\\
&\quad\Bigl[\int_{-\infty}^{\infty}(1+\alpha^{2})\,
\bigl|F[f_{n}](\alpha)-\mathcal F[f](\alpha)\bigr|^{2}\,d\alpha\Bigr]^{1/2}=B_{n},\quad\text{say},
\end{aligned}
\]
where $B_{n}\to 0$, as $n\to\infty$. Hence
\begin{equation}\tag{5.7}\label{eq:II-5.7}
\mathcal F[f_{n}]-\mathcal F[f]\in L_{1}(-\infty,\infty),\qquad\text{for }n=1,2,\dots.
\end{equation}
But $f_{n}-f\in L_{2}(-\infty,\infty)$, so that
\[
\begin{aligned}
&f_{n}(x) - f(x) = \lim_{R\to\infty\atop(L_{2}\text{-norm})} \frac{1}{\sqrt{2\pi}}  \int_{-R}^{R}\\
&\quad \bigl[F[f_{n}](\alpha)-F[f](\alpha)\bigr]e^{-i\alpha x}\,d\alpha\\
&= \frac{1}{\sqrt{2\pi}}  \int_{-\infty}^{\infty}  \bigl[F[f_{n}](\alpha)-F[f](\alpha)\bigr]e^{-i\alpha x}\,d\alpha,
\end{aligned}
\]
almost everywhere, by \eqref{eq:II-5.7}. Hence we have, almost everywhere,
\begin{equation}\tag{5.7'}\label{eq:II-5.7p}
|f_{n}(x)-f(x)|\le \|F[f_{n}]-F[f]\|_{1}\le B_{n}\to 0,\qquad\text{as }n\to\infty.
\end{equation}
Therefore there exists $c>0$, independent of $x$, such that
\begin{equation}\tag{5.8}\label{eq:II-5.8}
|f_{n}(x)-f(x)|<c<\infty,\qquad\text{almost everywhere}.
\end{equation}
Since $f^{*}\in L_{2}(-\infty,\infty)$ by definition, we have $f^{*}\in L_{1}(-R,R)$ for $0<R<\infty$. And
\[
\Bigl|\int_{-R}^{R}\bigl(f_{n}'\overline{f}_{n}-\overline{f^{*}}\overline{f}\bigr)(x)\,dx\Bigr|
\le\int_{-R}^{R}|f_{n}'|\,|\overline{f}_{n}-\overline{f}|\,dx
+\int_{-R}^{R}|\overline{f}|\,|f_{n}'-f^{*}|\,dx
\]
\begin{equation}\tag{5.9}\label{eq:II-5.9}
\le B_{n}\int_{-R}^{R}|f_{n}'|\,dx
+\Bigl(\int_{-R}^{R}|f|^{2}\,dx\Bigr)^{1/2}
\Bigl(\int_{-R}^{R}|f_{n}'-f^{*}|^{2}\,dx\Bigr)^{1/2}.
\end{equation}
Since $f\in L_{2}(-\infty,\infty)$, and $\|f_{n}'-f^{*}\|_{2}\to 0$, as $n\to\infty$, by \eqref{eq:II-5.6}, the second term on the right-hand side of \eqref{eq:II-5.9} tends to zero as $n\to\infty$. The first term on the right-hand side is
\begin{equation}\tag{5.10}\label{eq:II-5.10}
\begin{aligned}
&B_{n}\int_{-R}^{R}|f_{n}'-f^{*}+f^{*}|\,dx
\le B_{n}\int_{-R}^{R}|f_{n}'-f^{*}|\,dx+B_{n}\int_{-R}^{R}|f^{*}|\,dx\\
&\to 0,\qquad\text{as }n\to\infty\ \text{(for fixed }R>0),
\end{aligned}
\end{equation}
since $\|f_{n}'-f^{*}\|_{2}\to 0$. From \eqref{eq:II-5.9} and \eqref{eq:II-5.10} we have
\begin{equation}\tag{5.11}\label{eq:II-5.11}
\lim_{n\to\infty}\int_{-R}^{R}f_{n}'\overline{f}_{n}\,dx
=\int_{-R}^{R}f^{*}\overline{f}\,dx,
\end{equation}
and similarly also
\[
\int_{-R}^{R}x f_{n}'(x)\overline{f_{n}(x)}\,dx\to\int_{-R}^{R}x f^{*}(x)\overline{f(x)}\,dx,\qquad\text{as }n\to\infty,
\]
and
\[
\int_{-R}^{R}x\overline{f_{n}'(x)}f_{n}(x)\,dx\to\int_{-R}^{R}x\overline{f^{*}(x)}f(x)\,dx,\qquad\text{as }n\to\infty.
\]
Hence
\begin{equation}\tag{5.12}\label{eq:II-5.12}
\lim_{R\to\infty}\lim_{n\to\infty}\int_{-R}^{R}x\bigl(f_{n}'\overline{f}_{n}+\overline{f_{n}'}f_{n}\bigr)\,dx
=\int_{-\infty}^{\infty}x\bigl(f^{*}\overline{f}+\overline{f^{*}}f\bigr)\,dx,
\end{equation}
where
\begin{equation}\tag{5.12'}\label{eq:II-5.12p}
\Bigl|\int_{-\infty}^{\infty}x\bigl(f^{*}\overline{f}+\overline{f^{*}}f\bigr)\,dx\Bigr|<\infty,
\end{equation}
by \eqref{eq:II-5.3} together with the hypotheses $x|f(x)|\in L_{2}(-\infty<x<\infty)$, $\alpha|F[f](\alpha)|\in L_{2}(-\infty<\alpha<\infty)$.

Now
\[
\int_{-R}^{R}x\bigl(f_{n}'(x)\overline{f_{n}(x)}+\overline{f_{n}'(x)}f_{n}(x)\bigr)\,dx
=\int_{-R}^{R}x\bigl(|f_{n}(x)|^{2}\bigr)'\,dx
\]
\begin{equation}\tag{5.13}\label{eq:II-5.13}
=\lim_{R\to\infty}\Bigl\{R\bigl[|f(R)|^{2}+|f(-R)|^{2}\bigr]\Bigr\}-\|f\|_{2}^{2},
\end{equation}
since $f_{n}\to f$ almost everywhere, as $n\to\infty$, by \eqref{eq:II-5.7p}, and $\|f_{n}-f\|_{2}\to 0$, as $n\to\infty$ [assuming, as we may, that $R\to\infty$ through values which are not in the exceptional null set in \eqref{eq:II-5.7p}]. The left-hand side of \eqref{eq:II-5.13} is finite by \eqref{eq:II-5.12} and \eqref{eq:II-5.12p}; so is $\|f\|_{2}$; hence
\[
\lim_{R\to\infty}R\Bigl\{|f(R)|^{2}+|f(-R)|^{2}\Bigr\}
\]
is finite, and $\ge 0$. If the limit is $\delta>0$, then there exists $R_{0}$ such that we have
\[
\int_{R_{0}}^{\infty}\Bigl\{|f(R)|^{2}+|f(-R)|^{2}\Bigr\}\,dR
>\frac{\delta}{2}\int_{R_{0}}^{\infty}\frac{dR}{R}=\infty,
\]
which contradicts the assumption $f\in L_{2}(-\infty,\infty)$. Hence
\[
\lim_{R\to\infty}R\Bigl\{|f(R)|^{2}+|f(-R)|^{2}\Bigr\}=0,
\]
and this, taken together with \eqref{eq:II-5.13} and \eqref{eq:II-5.12}, yields
\[
\Bigl[\int_{-\infty}^{\infty}x\bigl(f^{*}(x)\overline{f}(x)+\overline{f^{*}}(x)f(x)\bigr)\,dx\Bigr]^{2}=\|f\|_{2}^{4},
\]
so that \eqref{eq:II-5.4} is proved, hence the inequality in the theorem.

In order to determine when the inequality becomes an equality, we note (by the first application of Schwarz's inequality just before \eqref{eq:II-5.3}) that
\[
\int_{-\infty}^{\infty}x^{2}\,|f(x)|^{2}\,dx
\int_{-\infty}^{\infty}\alpha^{2}\,|F[f](\alpha)|^{2}\,d\alpha
=\Bigl[\int_{-\infty}^{\infty}\bigl|x f(x)f^{*}(x)\bigr|^{2}\,dx\Bigr]^{2}
\]
if (and only if) $f^{*}(x)=K x\,\overline{f}(x)$ almost everywhere, $K$ being a complex constant. Here $f^{*}$ is defined as in \eqref{eq:II-5.2}. In fact, $f^{*}$ is the derivative almost everywhere of $f$. For by \eqref{eq:II-5.6}, we have
\[
\|f_{n}'-f^{*}\|_{2}\to 0,\qquad\text{as }n\to\infty,
\]
which implies that
\[
\lim_{n\to\infty}\int_{0}^{x}f_{n}'(y)\,dy=\int_{0}^{x}f^{*}(y)\,dy
\]
over any finite interval $[0,x]$. But the left-hand side equals
\[
\lim_{n\to\infty}[f_{n}(x)-f_{n}(0)]=f(x)-f(0),
\]
almost everywhere, by \eqref{eq:II-5.7p}.

Hence $f$ equals, almost everywhere, an absolutely continuous function, and
\begin{equation}\tag{5.14}\label{eq:II-5.14}
f^{*}(x)=\frac{d}{dx}\,[f(x)],\qquad\text{for almost all }x\in(-\infty,\infty).
\end{equation}
With this identification, \eqref{eq:II-5.1} becomes an equality if
\[
\frac{d}{dx}\,(f(x))=K x\,\overline{f}(x),\qquad(K\text{ a complex constant}),
\]
or $\frac{1}{x}\frac{d}{dx}f(x)=K\overline{f}(x)$, or $\frac{1}{x}\frac{d}{dx}\bigl(\frac{1}{x}\frac{d}{dx}f(x)\bigr)=\frac{1}{x}\frac{d}{dx}(K\overline{f})=|K|^{2}f(x)$, that is to say, $\bigl(\frac{1}{x}\frac{d}{dx}\bigr)^{2}f(x)-|K|^{2}f(x)=0$, or $(D^{2}-|K|^{2})f=0$, where $D=\frac{1}{x}\frac{d}{dx}$, which implies that $(D-|K|)(D+|K|)f=0$.
If $(D+|K|)f=0$, then $\frac{1}{x}\frac{df}{dx}=-|K|f$, or $\frac{f'}{f}=-|K|x$, or $\log f=-|K|\frac{x^{2}}{2}+c$,
or
\begin{equation}\tag{5.15}\label{eq:II-5.15}
f(x)=c_{1}\,e^{-|K|x^{2}/2}\in L_{2}(-\infty,\infty).
\end{equation}
If $(D-|K|)f=0$, then $f(x)=c_{2}e^{+|K|x^{2}/2}\notin L_{2}(-\infty,\infty)$.

The function $f(x)=e^{-|\kappa|x^{2}/2}$ actually satisfies the equality in \eqref{eq:II-5.1}, for if $f^{*}(x)=K x\overline{f}(x)$, then we have
\[
\begin{aligned}
&\int_{-\infty}^{\infty}x^{2}\,|f(x)|^{2}\,dx
\int_{-\infty}^{\infty}\alpha^{2}\,|F[f](\alpha)|^{2}\,d\alpha = \Bigl(\int_{-\infty}^{\infty}|x f^{*}(x)\overline{f}(x)|\,dx\Bigr)^{2}\\
&= \Bigl(\int_{-\infty}^{\infty}|K|x^{2}|f(x)|^{2}\,dx\Bigr)^{2}\\
&= |K|^{2}\Bigl(\int_{-\infty}^{\infty}x^{2}|f(x)|^{2}\,dx\Bigr)^{2},
\end{aligned}
\]
and if $f(x)=e^{-|\kappa|x^{2}/2}$, the last expression is
\[
\begin{aligned}
&|K|^{2}\Bigl(\int_{-\infty}^{\infty}x^{2}e^{-|K|x^{2}}\,dx\Bigr)^{2} = |K|^{2}\Bigl(\int_{-\infty}^{\infty}x\frac{-2|K|x}{-2|K|}e^{-|K|x^{2}}\,dx\Bigr)^{2}\\
&= |K|^{2}\Bigl(\int_{-\infty}^{\infty}\frac{e^{-|K|x^{2}}}{2|K|}\,dx\Bigr)^{2}\\
&= |K|^{2}\Bigl(\frac{1}{2|K|^{3/2}}\int_{-\infty}^{\infty}e^{-t^{2}}\,dt\Bigr)^{2}\\
&= |K|^{2}\frac{1}{4|K|^{3}}(\sqrt{\pi})^{2}\\
&= \frac{\pi}{4|K|},
\end{aligned}
\]
while
\[
\begin{aligned}
&\frac{1}{4}\|f\|_{2}^{4} = \frac{1}{4}\Bigl(\int_{-\infty}^{\infty}|f(x)|^{2}\,dx\Bigr)^{2}\\
&= \frac{1}{4}\Bigl(\int_{-\infty}^{\infty}e^{-|K|x^{2}}\,dx\Bigr)^{2}\\
&= \frac{1}{4}\Bigl(\frac{1}{\sqrt{|K|}}\int_{-\infty}^{\infty}e^{-t^{2}}\,dt\Bigr)^{2}\\
&= \frac{1}{4|K|}(\sqrt{\pi})^{2}\\
&= \frac{\pi}{4|K|}.
\end{aligned}
\]
Hence the equality holds, and the theorem is proved.
\end{proof}

\setcounter{section}{5}
\section{Hardy's theorem}
\label{sec:II-6}

We have noted that the function $f(x)=e^{-x^{2}}$ has the special property that
\[
\int_{-\infty}^{\infty}e^{-x^{2}}e^{i\alpha x}\,dx=e^{i\pi}e^{-\alpha^{2}/4},\qquad-\infty<\alpha<\infty,
\]
(cf.\ \S\ref{sec:I-1}, Ch.~I), which implies that
\[
\frac{1}{\sqrt{2\pi}}\int_{-\infty}^{\infty}e^{-x^{2}/2}e^{i\alpha x}\,dx=e^{-\alpha^{2}/2}.
\]
Hardy has shown that this property of the function $e^{-x^{2}/2}$, together with its order of magnitude, characterize the function in the sense of the following

\begin{theorem}\label{thm:II-9}
Let $f(x)$ be a measurable function defined on $-\infty<x<\infty$. Let $f(x)=O(e^{-x^{2}/2})$, as $|x|\to\infty$, and
\[
F[f](\alpha)=\frac{1}{\sqrt{2\pi}}\int_{-\infty}^{\infty}f(x)e^{i\alpha x}\,dx
=\mathcal O(e^{-\alpha^{2}/2}),
\]
as $|\alpha|\to\infty$. Then
\[
f(x)=c\,e^{-x^{2}/2},
\]
where $c$ is a complex number.
\end{theorem}

For the proof we need to apply the principle of Phragm\'en--Lindel\"of, which may be viewed as an extension of the maximum principle.

\begin{nstmt}{6.1}\label{eq:II-6.1}
\textbf{Theorem of Phragm\'en--Lindel\"of}. Let $f(z)$ be a non-constant holomorphic function of the complex variable $z$ ($=re^{i\theta}$), in the domain $D$ defined by the relations $r\ge 0$, $-\frac{\pi}{2\alpha}\le\theta\le\frac{\pi}{2\alpha}$, $\alpha>\frac{1}{2}$, and let $|f(z)|\le M<\infty$, for $r\ge 0$, $\theta=\pm\frac{\pi}{2\alpha}$. If $|f(z)|<\kappa e^{r^{\beta}}$, $\beta<\alpha$, $z\in\overline{D}$, where $\overline{D}$ denotes the closure of $D$, and the constant $\kappa$ is independent of $z$, then we have $|f(z)|<M$, $z\in D$.
\end{nstmt}
\begin{proof}
Consider the function
\[
F(z)=e^{-\varepsilon z^{\gamma}}f(z),\qquad \beta<\gamma<\alpha,\quad \varepsilon>0.
\]
We have
\[
|F(z)|=e^{-\varepsilon r^{\gamma}\cos\gamma\theta}|f(z)|.
\]
On the lines $\theta=\pm\frac{\pi}{2\alpha}$, we have $\cos\gamma\theta>0$, since $\gamma<\alpha$. Hence on these lines
\[
|F(z)|\le |f(z)|\le M.
\]
Further on the arc defined by $|z|=R>0$, $|\theta|\le\frac{\pi}{2\alpha}$, we have
\[
\begin{aligned}
&|F(z)| \le e^{\varepsilon R^{\gamma}\cos(\frac{1}{2}\gamma\frac{\pi}{\alpha})}|f(z)| < \kappa\cdot e^{R^{\beta}-\varepsilon R^{\gamma}\cos(\frac{1}{2}\frac{\gamma\pi}{\alpha})}\\
&\to 0,  \qquad\text{as }R \to \infty,
\end{aligned}
\]
since $\beta<\gamma<\alpha$. Hence for $R\ge R_{0}>0$, we have $|F(z)|\le M$, on that arc. By the maximum principle, we have
\[
|F(z)|\le M,\qquad 0\le r\le R,\quad |\theta|\le\frac{\pi}{2\alpha}.
\]
It follows that
\[
|f(z)|\le M\cdot e^{\varepsilon r^{\gamma}},\qquad z\in D,\quad \varepsilon>0.
\]
On letting $\varepsilon\downarrow 0$, we get the required result.
\end{proof}

\begin{equation*}\tag{6.2}\label{eq:II-6.2}
\text{\textbf{Corollary}. Take $\alpha=2$, $\beta=1$.}
\end{equation*}

\begin{proof}[Proof of Theorem~\ref{thm:II-9}]
Let $z=x+iy$, where $x$ and $y$ are real, $i=\sqrt{-1}$. Consider the function defined by
\begin{equation}\tag{6.3}\label{eq:II-6.3}
\tilde f(z)=\frac{1}{\sqrt{2\pi}}\int_{-\infty}^{\infty}f(t)e^{izt}\,dt.
\end{equation}
The integral converges absolutely and uniformly in every strip $-\infty<-A\le\operatorname{Im}z\le A<\infty$, since
\[
\begin{aligned}
&\frac{1}{\sqrt{2\pi}}\int_{-\infty}^{\infty}\bigl|f(t)e^{izt}\bigr|\,dt \le \kappa_{1}\int_{-\infty}^{\infty}e^{-\frac{1}{2}t^{2}-yt}\,dt\\
&= \kappa_{1}\int_{-\infty}^{\infty}e^{-\frac{1}{2}(t+y)^{2}+y^{2}/2}\,dt\\
&= \kappa_{2}\,e^{y^{2}/2},\quad\text{say}.
\end{aligned}
\]
Hence $\tilde f(z)$ is an entire function of $z$.

Since
\[
f(x)=\frac{1}{2}\bigl(f(x)+f(-x)\bigr)+\frac{1}{2}\bigl(f(x)-f(-x)\bigr)
=f_{1}(x)+f_{2}(x),
\]
say, where $f_{1}$ is an even function, and $f_{2}$ is an odd function, and $f_{1},f_{2}$ satisfy the same conditions as $f$, we shall consider the case $f$ even, and $f$ odd, separately.

Let $f$ be even. Then $\tilde f$ is also even, so that
\[
\tilde f(z)=\sum_{n=0}^{\infty}a_{n}z^{2n},
\]
and
\begin{equation}\tag{6.4}\label{eq:II-6.4}
\varphi(z):=\tilde f(\sqrt{z})=\sum_{n=0}^{\infty}a_{n}z^{n}
\end{equation}
is an entire function.

Set $z=r e^{i\theta}$, so that $x=r\cos\theta$, $y=r\sin\theta$, $\sqrt{z}=\sqrt{r}\,e^{i\theta/2}$, and $\operatorname{Im}(\sqrt{z})=r^{1/2}\sin(\theta/2)$. From \eqref{eq:II-6.4} we have
\begin{equation}\tag{6.5}\label{eq:II-6.5}
|\varphi(z)|=|\tilde f(\sqrt{z})|\le\kappa_{2}e^{\frac{1}{2}r\sin^{2}(\theta/2)}
\le\kappa_{2}e^{\frac{1}{2}r},
\end{equation}
for all $z$.

For $z$ real, and positive, say $z=r>0$, we have
\begin{equation}\tag{6.6}\label{eq:II-6.6}
|\varphi(r)|=|\tilde f(\sqrt{r})|\le\kappa_{3}e^{-\frac{1}{2}r},
\end{equation}
by hypothesis. Choose $M>\max(\kappa_{2},\kappa_{3})$, and $\alpha$ such that $0<\alpha<\pi$, and define the function
\[
\begin{aligned}
\omega(z,\alpha) &= \exp\Bigl[\frac{iz}{2}\frac{e^{-i\alpha/2}}{\sin(\alpha/2)}\Bigr]\\
&= \exp\Bigl[\frac{ir}{2}\frac{e^{i(\theta-\alpha/2)}}{\sin(\alpha/2)}\Bigr]\\
&= \exp\Bigl[\frac{ir}{2\sin(\alpha/2)}\bigl\{\cos(\theta-\tfrac{\alpha}{2})+i\sin(\theta-\tfrac{\alpha}{2})\bigr\}\Bigr]\\
&= \exp\Bigl[-\frac{r}{2}\frac{\sin(\theta-\alpha/2)}{\sin(\alpha/2)}+i\frac{r}{2}\frac{\cos(\theta-\alpha/2)}{\sin(\alpha/2)}\Bigr].
\end{aligned}
\]
Now
\[
|\omega(r,0,\alpha)|=e^{\frac{1}{2}r},\qquad |\omega(r,\alpha,\alpha)|=e^{-\frac{1}{2}r}.
\]
Hence (on writing $\varphi(r,\theta)\equiv\varphi(z)$, $\varphi(r)=\varphi(r,0)$), by \eqref{eq:II-6.6}, we have
\[
|\omega(r,0,\alpha)\cdot\varphi(r,0)|\le M,
\]
and
\[
|\omega(r,\alpha,\alpha)\cdot\varphi(r,\alpha)|\le M,\qquad\text{by \eqref{eq:II-6.5}}.
\]
By the Phragm\'en--Lindel\"of principle, we obtain
\[
|\omega(z)\,\varphi(z)|\le M,\qquad\text{for }0\le\theta\le\alpha,\ 0<\alpha<\pi;
\]
that is to say
\[
|\varphi(z)|\le M\exp\Bigl[\frac{r}{2}\cdot\frac{\sin(\theta-\alpha/2)}{\sin(\alpha/2)}\Bigr].
\]
If we keep $\theta$ fixed, and let $\alpha\uparrow\pi$, then
\[
|\varphi(z)|\le M\,e^{-\frac{1}{2}r\cos\theta},\qquad\text{for }0\le\theta<\pi.
\]
By continuity this holds also for $\theta=\pi$.

Similarly we consider the half-plane $-\pi<\theta<0$, and obtain
\[
\bigl|e^{\frac{1}{2}z}\varphi(z)\bigr|\le M,\qquad\text{for all (finite) }z.
\]
Since $\varphi$ is entire, $e^{\frac{1}{2}z}\varphi(z)$ reduces to a constant, therefore
\[
\varphi(z)=c_{1}\,e^{-\frac{1}{2}z},\qquad\text{or}\qquad \tilde f(z)=c\,e^{-\frac{1}{2}z^{2}},
\]
and the theorem follows.

If $f$ is odd, then $\tilde f$ is odd, hence $\tilde f(0)=0$, so that $\frac{\tilde f(z)}{z}$ is an even, entire function. By what we have just proved, $\tilde f(z)=c_{2}z e^{-\frac{1}{2}z^{2}}$. But for $z=x$ real, $\tilde f(x)=O(e^{-\frac{1}{2}x^{2}})$, by hypothesis. This is possible only if $c_{2}=0$ or $\tilde f(z)=f(z)\equiv 0$.
\end{proof}

\subsection*{The theorem of Paley and Wiener}
\noindent A fundamental theorem due to Paley and Wiener enables us to give a characterization of the Fourier transforms of functions belonging to $L_{2}(-\infty,\infty)$ which vanish outside a finite interval, in terms of entire functions of exponential type in the complex plane.

An entire function $f(z)$ of the complex variable $z$ is said to be of exponential type, if
\begin{equation}\tag{7.1}\label{eq:II-7.1}
f(z)=\mathcal O(e^{A|z|}),\qquad\text{as }z\to\infty,
\end{equation}
for some (finite) $A>0$. The lower bound of such numbers $A$ is called the type of $f$; it is non-negative.

We denote by $E^{\sigma}$, $\sigma\ge 0$, the class of entire functions of type at most $\sigma$. Thus if $f\in E^{\sigma}$, then
\begin{equation}\tag{7.2}\label{eq:II-7.2}
f(z)=O\bigl(e^{(\sigma+\varepsilon)|z|}\bigr),\qquad\text{as }z\to\infty,
\end{equation}
for every $\varepsilon>0$.

\begin{theorem}[Paley and Wiener]\label{thm:II-10}
Let $0<A<\infty$. Then we have
\begin{equation}\tag{7.3}\label{eq:II-7.3}
F(x)=\frac{1}{\sqrt{2\pi}}\int_{-A}^{A}f(u)e^{ixu}\,du,\qquad-\infty<x<\infty,
\end{equation}
for some $f\in L_{2}(-A,A)$, if and only if $F(x)\in L_{2}(-\infty<x<\infty)$ and $F$ can be extended to the complex plane as a function of the class $E^{A}$.
\end{theorem}
\begin{proof}
(First part.) If \eqref{eq:II-7.3} holds with $f(u)\in L_{2}(-A<u<A)$, then $f(u)\in L_{1}(-A<u<A)$, and for complex $z$,
\[
F(z)=\frac{1}{\sqrt{2\pi}}\int_{-A}^{A}f(u)e^{izu}\,du\qquad(z=x+iy)
\]
is an entire function of $z$, with
\[
|F(z)|\le\frac{1}{\sqrt{2\pi}}e^{A|y|}\int_{-A}^{A}|f(u)|\,du=O(e^{A|z|}),
\]
so that $F\in E^{A}$. Further the Fourier transform of $f$, where $f(u)$ is defined to be zero for $|u|>A$, is $F$, so that by Plancherel's theorem, $\|F\|_{2}=\|f\|_{2}<\infty$.

(Second part.) Let $F(x)\in L_{2}(-\infty<x<\infty)$, and for complex $z$, let $F(z)\in E^{A}$. Then by Plancherel's theorem,
\begin{equation}\tag{7.4}\label{eq:II-7.4}
f(x)=\lim_{R\to\infty\atop(L_{2}\text{-norm})}\frac{1}{\sqrt{2\pi}}\int_{-R}^{R}F(u)e^{-ixu}\,du
\end{equation}
belongs to $L_{2}(-\infty<x<\infty)$. We shall prove that $f(x)$ vanishes almost everywhere for $|x|>A>0$, so that $f(x)\in L_{2}(-A<x<A)$.

For complex $z$, let
\begin{equation}\tag{7.5}\label{eq:II-7.5}
g(z)=\int_{-1/2}^{1/2}F(u-z)\,du.
\end{equation}
Then
\begin{equation}\tag{7.6}\label{eq:II-7.6}
g(z)\text{ is an entire function of }z,
\end{equation}
such that, for $\varepsilon>0$,
\[
\begin{aligned}
&|g(z)| \le \frac{1}{2}  \int_{-1/2}^{1/2}|F(u-z)|\,du = O  \Bigl(\int_{-1/2}^{1/2}e^{(A+\varepsilon)(|z|+|u|)}\,du\Bigr)\\
&= O  \bigl(e^{(A+\varepsilon)|z|}\bigr),
\end{aligned}
\]
since $F\in E^{A}$. Further, for real $x$, we have
\begin{equation}\tag{7.7}\label{eq:II-7.7}
|g(x)|^{2}\le\int_{-1/2}^{1/2}|F(u-x)|^{2}\,du\le \|F\|_{2}^{2}<\infty,
\end{equation}
and
\[
\begin{aligned}
& \int_{-\infty}^{\infty}|g(x)|^{2}\,dx \le \int_{-\infty}^{\infty}  \Bigl\{  \int_{-1/2}^{1/2}|F(u-x)|^{2}\,du\Bigr\}dx\\
&= \int_{-1/2}^{1/2}du  \int_{-\infty}^{\infty}|F(u-x)|^{2}\,dx\\
&= \int_{-1/2}^{1/2}\|F\|_{2}^{2}\,du\\
&= \|F\|_{2}^{2}<\infty,
\end{aligned}
\]
so that $g(x)$ is bounded, with $g(x)\in L_{2}(-\infty<x<\infty)$.

Next let
\begin{equation}\tag{7.8}\label{eq:II-7.8}
G(z)=e^{iBz}g(z),\qquad z\text{ complex},\quad B>A>0.
\end{equation}
Then $G(z)$ is an entire function of exponential type, by \eqref{eq:II-7.6} and \eqref{eq:II-7.7}. Further if $z$ is real, and $z=x$, then
\begin{equation}\tag{7.9}\label{eq:II-7.9}
|G(x)|=|g(x)|=o(1),\qquad\text{by \eqref{eq:II-7.8}},
\end{equation}
and if $z=ib$, $b>0$, then by \eqref{eq:II-7.7}, we have
\[
|G(ib)|=|e^{-Bb}g(ib)|=O(e^{-Bb+(A+\varepsilon)b})=O(e^{-b(B-A-\varepsilon)})\to 0,
\]
if $\varepsilon$ is chosen sufficiently small, since $B>A$. In particular,
\begin{equation}\tag{7.10}\label{eq:II-7.10}
|G(ib)|=O(1),\qquad\text{as }b\to+\infty\text{ and for }B>A.
\end{equation}
Because of \eqref{eq:II-7.9} and \eqref{eq:II-7.10} and the Phragm\'en--Lindel\"of principle \eqref{eq:II-6.1}, it follows that $|e^{iBz}g(z)|=O(1)$, for $z=Re^{i\theta}$, $0\le\theta\le\frac{\pi}{2}$, or $|g(Re^{i\theta})|<c_{1}e^{BR\sin\theta}$, $0\le\theta\le\frac{\pi}{2}$. Similarly $|g(Re^{i\theta})|<c_{2}e^{BR\sin\theta}$, for $\frac{\pi}{2}<\theta\le\pi$. Hence we have
\begin{equation}\tag{7.11}\label{eq:II-7.11}
|g(Re^{i\theta})|=O(e^{BR\sin\theta}),\qquad R>0,\quad 0\le\theta\le\pi.
\end{equation}
Let $L>0$, and $x>B>A>0$, and consider the integral
\[
\int_{-R}^{R}\frac{e^{ixu}}{1-Liu}\,g(u)\,du,
\]
where $R$ is large enough to ensure that $LR>1$. By Cauchy's theorem applied to the semicircular contour defined by $|x|\le R$, $\operatorname{Im}z=0$; and $\operatorname{Im}z>0$, $|z|=R$ (see Fig.), we have
\[
\int_{-R}^{R}\frac{e^{ixu}}{1-Liu}\,g(u)\,du
=-i\int_{0}^{\pi}\frac{e^{ixRe^{i\theta}}}{1-LiRe^{i\theta}}\,g(Re^{i\theta})\,Re^{i\theta}\,d\theta
\]
so that
\begin{equation}\tag{7.12}\label{eq:II-7.12}
\begin{aligned}
&\Bigl|  \int_{-R}^{R}\frac{e^{ixu}g(u)}{1-Liu}\,du\Bigr|\\
&\le \frac{c\cdot R}{LR-1}  \int_{0}^{\pi}e^{-xR\sin\theta+B R\sin\theta}\,d\theta \to 0,  \qquad\text{as }R \to \infty,
\end{aligned}
\end{equation}
since $LR>1$, $x>B$, by choice, $\sin\theta>0$ for $0<\theta<\pi$, and
\begin{equation}\tag{7.13}\label{eq:II-7.13}
\int_{0}^{\pi}e^{-aR\sin\theta}\,d\theta=O\Bigl(\frac{1}{R}\Bigr),\qquad\text{for any }a>0,
\end{equation}
which can be seen as follows. We have
\[
\int_{0}^{\pi}e^{-aR\sin\theta}\,d\theta
=\int_{0}^{\pi/2}+\int_{\pi/2}^{\pi}=I_{1}+I_{2},\quad\text{say},
\]
where
\[
|I_{1}|\le\Bigl|\int_{0}^{\pi/2}e^{-cR\theta}\,d\theta\Bigr|
=\Bigl|\frac{-e^{-cR\pi/2}}{cR}+\frac{1}{cR}\Bigr|
=O\Bigl(\frac{1}{R}\Bigr),\qquad c=\frac{2a}{\pi}>0,
\]
since $\frac{2}{\pi}\le\frac{\sin\theta}{\theta}\le 1$, for $0\le\theta\le\frac{\pi}{2}$, so that $-aR\sin\theta\le\frac{-2a}{\pi}R\theta=-cR\theta$. And
\[
\begin{aligned}
&\int_{\pi/2}^{\pi}e^{-aR\sin\theta}\,d\theta\\
&= -  \int_{\pi/2}^{0}e^{-aR\sin\varphi}\,d\varphi = I_{1},  \qquad(\text{with }\theta=\pi-\varphi,\ \sin\theta=\sin\varphi).
\end{aligned}
\]
Let
\begin{equation}\tag{7.14}\label{eq:II-7.14}
g_{-L}(u)=\frac{g(u)}{1-Liu},\qquad L>0,\quad -\infty<u<\infty.
\end{equation}
Then $g_{-L}\in L_{1}(-\infty,\infty)$, since it is the product of two functions each of which belongs to $L_{2}(-\infty,\infty)$. From \eqref{eq:II-7.12} we deduce that the Fourier transform of $g_{-L}$ vanishes for $L>0$, $x>B$;
\begin{equation}\tag{7.15}\label{eq:II-7.15}
F[g_{-L}](x)=0,\qquad L>0,\quad x>B.
\end{equation}
We have, however,
\begin{equation}\tag{7.16}\label{eq:II-7.16}
\int_{-\infty}^{\infty}\Bigl|\frac{g(u)}{1-Liu}-g(u)\Bigr|^{2}\,du\to 0,\qquad\text{as }L\downarrow 0,
\end{equation}
by Lebesgue's theorem on dominated convergence (cf.\ \eqref{eq:II-7.8}). Hence, by Plancherel's theorem,
\begin{equation}\tag{7.17}\label{eq:II-7.17}
\int_{-\infty}^{\infty}\bigl|F[g_{-L}](u)-F[g](u)\bigr|^{2}\,du\to 0,\qquad\text{as }L\downarrow 0.
\end{equation}
From \eqref{eq:II-7.15} it follows that
\begin{equation}\tag{7.18}\label{eq:II-7.18}
F[g](x)=0,\qquad\text{almost everywhere, for }x>B.
\end{equation}
We have, however, for real $v$,
\begin{equation}\tag{7.19}\label{eq:II-7.19}
g(v)=\int_{-1/2}^{1/2}F(u-v)\,du
=\frac{1}{\sqrt{2\pi}}\int_{-\infty}^{\infty}f(x)\,\frac{\sin x/2}{x/2}\,e^{-ivx}\,dx,
\end{equation}
by \eqref{eq:II-2.26}. (Note that $\check F[f](x)=F[f](-x)$.) Since $f(x)\frac{\sin x/2}{x/2}\in L_{1}(-\infty,\infty)\cap L_{2}(-\infty,\infty)$, we deduce from \eqref{eq:II-7.19} that
\[
F[g](x)=f(x)\,\frac{\sin x/2}{x/2},
\]
for almost all $x\in(-\infty,\infty)$. Since \eqref{eq:II-7.18} holds for every $B>A$, we conclude that $f(x)=0$ for almost every $x>A$.

Similarly we show that $f(x)=0$ for almost every $x<-A$. For we have only to consider $G(-ib)$, $b>0$, instead of $G(ib)$, and note that, after \eqref{eq:II-7.7},
\begin{equation}\tag{7.20}\label{eq:II-7.20}
|G(-ib)|=|e^{Bb}g(-ib)|=O(e^{b(A+\varepsilon+B)}),\qquad b>0,
\end{equation}
for every $\varepsilon>0$, and if $B<-A<0$, then for $\varepsilon$ sufficiently small, we have
\begin{equation}\tag{7.21}\label{eq:II-7.21}
|G(-ib)|=O(1),\qquad\text{as }b\to+\infty,\quad B<-A<0,
\end{equation}
corresponding to \eqref{eq:II-7.10}. Again by the Phragm\'en--Lindel\"of principle applied to a semi-circular domain in the lower half-plane, we obtain
\begin{equation}\tag{7.22}\label{eq:II-7.22}
|g(Re^{i\theta})|=O(e^{BR\sin\theta}),\qquad -\pi\le\theta\le 0,\quad R>0,
\end{equation}
corresponding to \eqref{eq:II-7.11}. And we have
\[
\begin{aligned}
&\Bigl|\int_{-R}^{R}\frac{e^{ixu}g(u)}{1+Liu}\,du\Bigr| \le \Bigl|i\int_{-\pi}^{0}\frac{e^{ixRe^{i\theta}}g(Re^{i\theta})\,Re^{i\theta}}{1+LiRe^{i\theta}}\,d\theta\Bigr|\\
&\le \frac{cR}{LR-1}\int_{-\pi}^{0}e^{-(x-B)R\sin\theta}\,d\theta\\
&= \frac{cR}{LR-1}\int_{0}^{\pi}e^{(x-B)R\sin\theta}\,d\theta,
\end{aligned}
\]
which tends to zero as $R\to\infty$ if $x<B$, $c$ being a constant. This corresponds to \eqref{eq:II-7.12}. And we deduce (as in \eqref{eq:II-7.15}) that
\begin{equation}\tag{7.23}\label{eq:II-7.23}
F[g_{L}](x)=0,\qquad\text{for }L>0,\quad x<B<-A<0,
\end{equation}
where $g_{L}(x)=\frac{g(x)}{1+Lix}\in L_{1}(-\infty<x<\infty)$. It then follows, as before, that $f(x)=0$ for almost every $x<-A$, and the second part of the theorem is proved.
\end{proof}

Combining the Paley--Wiener theorem with the Riemann--Lebesgue theorem (Ch.~I, Th.~\ref{thm:I-1}), we obtain the following

\begin{nstmt}{7.24}\label{eq:II-7.24}
\textbf{Corollary}. If $F(z)$ is an entire function of exponential type, with $z=x+iy$, and $F(x)\in L_{2}(-\infty<x<\infty)$, then $F(x)\to 0$ as $|x|\to\infty$.
\end{nstmt}

\subsection*{Fourier series in $L_{2}(a,b)$}
\setcounter{section}{7}
\section{Fourier series in \texorpdfstring{$L_2(a,b)$}{L2(a,b)}}
\label{sec:II-8}

\noindent Let $\varphi_{0}(x)$, $\varphi_{1}(x)$, $\varphi_{2}(x)$, \dots, be complex-valued, non-null functions (that is, functions which are not almost everywhere equal to zero) defined on the interval $(a,b)$ of the real line. We say that $\{\varphi_{n}\}$ is an orthogonal set if
\[
\int_{a}^{b}\varphi_{m}(x)\overline{\varphi}_{n}(x)\,dx=
\begin{cases}
0, & m\neq n;\ m,n=0,1,2,\dots,\\
\lambda_{m}>0, & m=n.
\end{cases}
\]
We note that $\varphi_{m}(x)\in L_{2}(a<x<b)$, and that no $\varphi_{m}(x)$ can vanish identically, since that would imply that $\lambda_{m}=0$.

If, in addition, $\lambda_{m}=1$ for $m=0,1,2,\dots$, we call $\{\varphi_{n}\}$ an orthonormal set.

If $\{\varphi_{n}\}$ is orthogonal, then clearly $\{\varphi_{n}/\lambda_{n}^{1/2}\}$ is orthonormal.

Given $f\in L_{2}(a,b)$, let
\[
c_{n}=\frac{1}{\lambda_{n}}\int_{a}^{b}f(x)\overline{\varphi}_{n}(x)\,dx,
\]
for any integer $n\ge 0$. We call $c_{n}$ the $n^{\text{th}}$ Fourier coefficient of $f$, and write
\begin{equation}\tag{8.1}\label{eq:II-8.1}
f\sim(c_{n}).
\end{equation}
We call the series
\[
\sum_{n=0}^{\infty}c_{n}\varphi_{n}(x)
\]
the Fourier (orthogonal) series of $f$, relative to $(\varphi_{n})$, and indicate that relation by writing
\begin{equation}\tag{8.2}\label{eq:II-8.2}
f(x)\sim\sum_{n=0}^{\infty}c_{n}\,\varphi_{n}(x).
\end{equation}
The set $\{e^{inx}\}$, $n=0,\pm1,\pm2,\dots$, is orthogonal on $(0,2\pi)$, while the set $\{e^{inx}(2\pi)^{-1/2}\}$ is orthonormal, since
\[
\int_{\alpha}^{\alpha+2\pi}e^{imx}\cdot e^{-inx}\,dx=
\begin{cases}
2\pi, & m=n,\\
0, & m\neq n,
\end{cases}
\]
for any real $\alpha$.

In order to have the equality
\[
f(x)=\sum_{n=0}^{\infty}c_{n}\,\varphi_{n}(x),\qquad\text{for any }f\in L_{2}(a,b),
\]
it is necessary that the set $\{\varphi_{n}\}$ should be ``complete'', in the sense that if we add a new, non-null function $\psi$ to the set, then the resulting set is no longer orthogonal. Otherwise there would exist a non-null function, namely $\psi$ itself, with all its Fourier coefficients equal to zero.

We thus define a set $(\varphi_{n})$, $\varphi_{n}\in L_{2}(a,b)$, to be complete, if there exists no non-null function in $L_{2}(a,b)$ which is orthogonal to $\varphi_{n}$ for every $n>0$. In other words,
\[
\int_{a}^{b}f(x)\overline{\varphi}_{n}(x)\,dx=0,\qquad n=0,1,2,\dots,
\]
implies that $f(x)=0$, for almost every $x\in(a,b)$.

It follows that if $\{\varphi_{n}\}$ is complete, and orthonormal, and the functions $f,g\in L_{2}(a,b)$ have the same Fourier coefficients, then $f$ is equivalent to $g$ (cf.\ \S\ref{sec:I-1}, Ch.~I), that is to say that $f$ equals $g$ almost everywhere in $(a,b)$. In this sense, the Fourier series of $f$ is unique.

Let $f\in L_{2}(a,b)$, with $\|f\|_{2}=(\int_{a}^{b}|f(x)|^{2}\,dx)^{1/2}$. Let $(\varphi_{n})$ be an orthonormal set in $L_{2}(a,b)$, and let $c_{n}$ denote the $n$th Fourier coefficient of $f$ relative to $(\varphi_{n})$. By a polynomial in the $\varphi_{n}$, of rank $k$, where $k$ is a non-negative integer, is meant an expression of the form
\begin{equation}\tag{8.3}\label{eq:II-8.3}
\Phi_{k}(x)=\sum_{m=0}^{k}\gamma_{m}\,\varphi_{m}(x),
\end{equation}
where all the $\gamma$'s are complex numbers. We call
\begin{equation}\tag{8.4}\label{eq:II-8.4}
f_{k}(x)=\sum_{m=0}^{k}c_{m}\,\varphi_{m}(x)
\end{equation}
the Fourier polynomial, of rank $k$, of the function $f$.

With this notation we have the following identity:
\begin{equation}\tag{8.5}\label{eq:II-8.5}
\|f-\Phi_{n}\|_{2}^{2}
=\|f\|_{2}^{2}-\sum_{m=0}^{n}|c_{m}|^{2}
+\sum_{m=0}^{n}|c_{m}-\gamma_{m}|^{2},
\end{equation}
and, in particular,
\begin{equation}\tag{8.6}\label{eq:II-8.6}
\|f-f_{n}\|_{2}^{2}=\|f\|_{2}^{2}-\sum_{m=0}^{n}|c_{m}|^{2}.
\end{equation}
To prove \eqref{eq:II-8.5} we note that
\[
\int_{a}^{b}f(x)\overline{\Phi}_{n}(x)\,dx
=\sum_{m=0}^{n}c_{m}\,\overline{\gamma}_{m},\qquad
\int_{a}^{b}|\Phi_{n}(x)|^{2}\,dx
=\sum_{m=0}^{n}|\gamma_{m}|^{2},
\]
and
\[
\begin{aligned}
\|f-\Phi_{n}\|_{2}^{2} &= \int_{a}^{b}\bigl\{f(x)-\Phi_{n}(x)\bigr\}\bigl\{\overline{f}(x)-\overline{\Phi}_{n}(x)\bigr\}\,dx\\
&= \|f\|_{2}^{2}-\sum_{m=0}^{n}c_{m}\overline{\gamma}_{m}
-\sum_{m=0}^{n}\overline{c}_{m}\gamma_{m}
+\sum_{m=0}^{n}|\gamma_{m}|^{2}.
\end{aligned}
\]
By \eqref{eq:II-8.6}, $\sum_{m=0}^{n}|c_{m}|^{2}\le\|f\|_{2}^{2}$, and on letting $n\to\infty$, we obtain Bessel's inequality, namely
\begin{equation}\tag{8.7}\label{eq:II-8.7}
\sum_{n=0}^{\infty}|c_{n}|^{2}\le\|f\|_{2}^{2}
=\int_{a}^{b}|f(x)|^{2}\,dx.
\end{equation}
From \eqref{eq:II-8.5} we conclude that
\begin{nstmt}{8.8}\label{eq:II-8.8}
the best approximation, in the $L_{2}$-norm, to $f\in L_{2}(a,b)$ by polynomials $\Phi_{n}$, of a given rank $n$, is provided by the Fourier polynomial $f_{n}$ of $f$.
\end{nstmt}

An important result on the Fourier series of functions in $L_{2}(a,b)$ is the following

\begin{nstmt}{8.9}\label{eq:II-8.9}
\textbf{The Riesz--Fischer theorem}. Given any sequence $(c_{n})$ of complex numbers, such that $\sum_{n=0}^{\infty}|c_{n}|^{2}<\infty$, there exists a function $f\in L_{2}(a,b)$, such that the $c_{n}$'s are the Fourier coefficients of $f$, relative to the given orthonormal set $(\varphi_{n})$. We have further:
\end{nstmt}
\begin{equation}\tag{8.10}\label{eq:II-8.10}
\int_{a}^{b}|f_{n}(x)-f(x)|^{2}\,dx\to 0,\qquad\text{as }n\to\infty,
\end{equation}
and
\begin{equation}\tag{8.11}\label{eq:II-8.11}
\int_{a}^{b}|f(x)|^{2}\,dx=\sum_{n=0}^{\infty}|c_{n}|^{2}.
\end{equation}
\begin{proof}
Let $f_{m}\equiv f_{m}(x)=\sum_{k=0}^{m}c_{k}\varphi_{k}(x)\equiv\sum_{k=0}^{m}c_{k}\varphi_{k}$. If $n>m$, then
\[
\int_{a}^{b}|f_{m}(x)-f_{n}(x)|^{2}\,dx=\sum_{k=m+1}^{n}|c_{k}|^{2}\to 0,\qquad\text{as }m,n\to\infty,
\]
by the hypothesis on the $c_{n}$. Since the function space $L_{2}(a,b)$ is complete (cf.\ \S\ref{sec:II-2}) (which is not to be confused with the set $(\varphi_{n})$ being complete, which is not assumed here), there exists a function $f\in L_{2}(a,b)$, such that
\begin{equation}\tag{8.12}\label{eq:II-8.12}
\int_{a}^{b}|f_{n}(x)-f(x)|^{2}\,dx\to 0,\qquad\text{as }n\to\infty.
\end{equation}
It follows that
\[
\int_{a}^{b}f(x)\overline{\varphi}_{m}(x)\,dx
=\lim_{n\to\infty}\int_{a}^{b}f_{n}(x)\overline{\varphi}_{m}(x)\,dx
=\lim_{n\to\infty}c_{m}=c_{m},
\]
so that $c_{m}$ is the $m^{\text{th}}$ Fourier coefficient of $f$, relative to $(\varphi_{n})$.

(Note that if $\|f-f_{n}\|_{2}\to 0$, and $\|g-g_{n}\|_{2}\to 0$, then $\|fg-f_{n}g_{n}\|_{1}\to 0$.) Hence, by \eqref{eq:II-8.5}, we have
\[
\|f-f_{n}\|_{2}^{2}=\|f\|_{2}^{2}-\sum_{m=0}^{n}|c_{m}|^{2},
\]
and \eqref{eq:II-8.12} now implies \eqref{eq:II-8.11}.
\end{proof}

\begin{remark}
We note that by \eqref{eq:II-8.6}, \eqref{eq:II-8.10} is equivalent to \eqref{eq:II-8.11}. We next make use of the completeness of $(\varphi_{n})$.
\end{remark}

\begin{nstmt}{8.13}\label{eq:II-8.13}
\textbf{Parseval's theorem}. (A) If $f\in L_{2}(a,b)$, and $(\varphi_{n})$ is a \textbf{complete}, orthonormal set in $L_{2}(a,b)$, and $c_{n}$ denotes the $n^{\text{th}}$ Fourier coefficient of $f$, relative to the $(\varphi_{n})$, then
\end{nstmt}
\begin{equation}\tag{8.14}\label{eq:II-8.14}
\int_{a}^{b}|f_{n}(x)-f(x)|^{2}\,dx\to 0,\qquad\text{as }n\to\infty,
\end{equation}
and
\begin{equation}\tag{8.15}\label{eq:II-8.15}
\int_{a}^{b}|f(x)|^{2}\,dx=\sum_{n=0}^{\infty}|c_{n}|^{2}.
\end{equation}
(B). If $f,g\in L_{2}(a,b)$, $f\sim(c_{n})$ and $g\sim(d_{n})$, (as defined in \eqref{eq:II-8.1}), then
\[
\int_{a}^{b}f(x)\overline{g}(x)\,dx
=\sum_{n=0}^{\infty}c_{n}\,\overline{d}_{n},
\]
the series on the right converging absolutely.
\begin{proof}
(A) By Bessel's inequality we have $\sum_{n=0}^{\infty}|c_{n}|^{2}<\infty$, and the Riesz--Fischer theorem gives an $f\in L_{2}(a,b)$ satisfying \eqref{eq:II-8.14} and \eqref{eq:II-8.15}. The completeness of $(\varphi_{n})$ implies that the $f$ differs from the given $f$ only on a set of measure zero in $(a,b)$, and \eqref{eq:II-8.14} and \eqref{eq:II-8.15} follow.

(B) We have
\[
\int_{a}^{b}f_{n}(x)\overline{g}(x)\,dx
=\sum_{m=0}^{n}c_{m}\int_{a}^{b}\varphi_{m}(x)\overline{g}(x)\,dx
=\sum_{m=0}^{n}c_{m}\overline{d}_{m}.
\]
Since $\|f-f_{n}\|_{2}\to 0$, as $n\to\infty$, by \eqref{eq:II-8.14} it follows that
\[
\int_{a}^{b}f_{n}(x)\overline{g}(x)\,dx\to\int_{a}^{b}f(x)\overline{g}(x)\,dx,\qquad\text{as }n\to\infty,
\]
and since $\sum_{n=0}^{\infty}|c_{n}|^{2}$, and $\sum_{n=0}^{\infty}|d_{n}|^{2}$, converge by \eqref{eq:II-8.15}, we have $\sum_{m=0}^{\infty}|c_{m}\overline{d}_{m}|<\infty$.
\end{proof}

The closure and completeness of $(\varphi_{n})$ in $L_{2}(a,b)$:
\begin{nstmt}{8.16}\label{eq:II-8.16}
A set of functions $(\varphi_{n})$, $\varphi_{n}\in L_{2}(a,b)$, is said to be \textbf{closed} in $L_{2}(a,b)$, if the polynomials in $\varphi_{n}$ (see \eqref{eq:II-8.3}) form a dense subset of $L_{2}(a,b)$ in the $L_{2}$-norm.
\end{nstmt}
\begin{nstmt}{8.17}\label{eq:II-8.17}
The orthonormal set $(\varphi_{n})$ is closed in $L_{2}(a,b)$, if and only if it is complete in $L_{2}(a,b)$.
\end{nstmt}
\begin{proof}
If $(\varphi_{n})$ is complete, then given any $f\in L_{2}(a,b)$, there exists, by \eqref{eq:II-8.14}, a sequence $(f_{n})$ of Fourier polynomials of $f$, such that $\|f-f_{n}\|_{2}\to 0$, as $n\to\infty$. Hence the set $(\varphi_{n})$ is closed.

Conversely let the set $(\varphi_{n})$ be closed in $L_{2}(a,b)$, and $f\in L_{2}(a,b)$, $f\sim(c_{n})$, and $c_{n}=0$ for all $n\ge 0$. Then there exists a sequence of polynomials in the $\varphi_{n}$, say $\Phi_{n}$, such that $\|\Phi_{n}-f\|_{2}\to 0$, as $n\to\infty$. By \eqref{eq:II-8.8}, we have $\|f_{n}-f\|_{2}\to 0$, as $n\to\infty$, where $f_{n}$ is the Fourier polynomial of $f$ of rank $n$. But $f_{n}=0$, since all the $c_{n}$'s are zero, hence $\|f\|_{2}=0$, or $f$ is zero almost everywhere in $(a,b)$. It follows that $(\varphi_{n})$ is complete in $L_{2}(a,b)$.
\end{proof}

\begin{remark}
1. Let $\varphi_{n}(x)=\frac{1}{\sqrt{2\pi}}e^{inx}$, $|x|\le\pi$, $n=0,\pm1,\pm2,\dots$. It is a classical result that $(\varphi_{n})$ is complete in $L_{1}(-\pi,\pi)$, hence also in $L_{2}(-\pi,\pi)$.

2. Let
\[
\varphi_{0}(x)=
\begin{cases}
+1, & 0<x<\frac{1}{2},\\
-1, & \frac{1}{2}<x<1,
\end{cases}
\]
and $\varphi_{0}(0)=\varphi_{0}(\frac{1}{2})=0$, with $\varphi(x)=\varphi(x+1)$. Let $\varphi_{n}(x)=\varphi_{0}(2^{n}x)$ for $n=0,1,2,\dots$. Then $(\varphi_{n})$ is orthonormal over $(0,1)$. The function $\varphi_{n}(x)$ takes alternately the values $+1$, and $-1$ in the intervals $(0,2^{-n-1})$, $(2^{-n-1},2\cdot2^{-n-1})$, $(2\cdot2^{-n-1},3\cdot2^{-n-1})$, \dots. If $m>n$, the integral of $\varphi_{m}\varphi_{n}$ over any of these intervals is zero, and the set $(\varphi_{n})$ is obviously orthonormal. It is not complete, since the function $\psi(x)=1$ may be added to it. (Note that the set $\{\psi\}\cup\{\varphi_{n}\}_{n>0}$ is not complete.) If we define
\[
\operatorname{sign}\alpha=1\quad\text{for }\alpha>0,\quad \operatorname{sign}\alpha=-1\quad\text{for }\alpha<0,\\
\text{and}\quad \operatorname{sign}\alpha=0\quad\text{for }\alpha=0,
\]
then we have $\varphi_{n}(x)=\operatorname{sign}\sin 2^{n+1}\pi x$. The functions $(\varphi_{n})$ are known as \emph{Rademacher's} functions.

3. The condition of orthonormality in \eqref{eq:II-8.17} is not necessary.
\end{remark}

\setcounter{section}{8}
\section{Hardy's interpolation formula}
\label{sec:II-9}

As an application of the Paley--Wiener theorem one can obtain some special formulae of interpolation at integer points for entire functions of exponential type, which help to determine the value of the function at an arbitrary point in the complex plane in terms of its values at the integer points on the real axis. The notions concerning Fourier series outlined in the previous section enter into the proofs.

If $F(z)$ is an entire function of $z$ belonging to the class $E^{\sigma}$, $\sigma>0$, as defined in \eqref{eq:II-7.2}, then $F(\frac{\pi z}{\sigma})\in E^{\pi}$, so that the study of functions in $E^{\sigma}$, $\sigma>0$, may be reduced to that of functions in $E^{\pi}$.

\begin{theorem}[Hardy]\label{thm:II-11}
Let $F(z)\in E^{\pi}$, $z=x+iy$, and
\begin{equation}\tag{9.1}\label{eq:II-9.1}
\int_{-\infty}^{\infty}|F(x)|^{2}\,dx<\infty.
\end{equation}
Then we have
\begin{equation}\tag{9.2}\label{eq:II-9.2}
F(z)=\frac{\sin\pi z}{\pi}\sum_{n=-\infty}^{\infty}(-1)^{n}\frac{F(n)}{z-n}.
\end{equation}
\end{theorem}
\begin{proof}
If we define
\[
\varphi_{n}(x)=
\begin{cases}
\frac{1}{\sqrt{2\pi}}e^{inx}, & |x|\le\pi,\\
0, & |x|>\pi,
\end{cases}
\]
for $n=0,\pm1,\pm2,\dots$, then the set $(\varphi_{n})$ is orthonormal over $(-\infty,\infty)$, but not complete on $(-\infty,\infty)$ (since any function which vanishes on $(-\pi,\pi)$ has all its Fourier coefficients equal to zero). It follows that the set $\{F[\varphi_{n}]\}$ is also orthonormal (by \eqref{eq:II-2.24}). Now
\[
\begin{aligned}
&\check F[\varphi_{n}](\alpha) = \frac{1}{\sqrt{2\pi}}\int_{-\infty}^{\infty}\varphi_{n}(x)e^{-i\alpha x}\,dx\\
&= \frac{1}{2\pi}\int_{-\pi}^{\pi}e^{i(n-\alpha)x}\,dx\\
&= \frac{\sin(\alpha-n)\pi}{\pi(\alpha-n)}\\
&= (-1)^{n}\frac{\sin(\alpha\pi)}{\pi(\alpha-n)}.
\end{aligned}
\]
Since $F(x)\in L_{2}(-\infty,\infty)$, let
\begin{equation}\tag{9.3}\label{eq:II-9.3}
F(x)\sim\sum_{n=-\infty}^{\infty}a_{n}\,\check F[\varphi_{n}](x)
=\sum_{n=-\infty}^{\infty}a_{n}(-1)^{n}\frac{\sin(\pi x)}{\pi(x-n)},
\end{equation}
where
\[
a_{n}=\frac{1}{\sqrt{2\pi}}\int_{-\infty}^{\infty}F(x)\,\overline{\check F[\varphi_{n}]}(x)\,dx
\]
is the $n^{\text{th}}$ Fourier coefficient of $F$ relative to $(\check F[\varphi_{n}])$. [Since $(\check F[\varphi_{n}])$ is not complete, the series in \eqref{eq:II-9.3} does not, in general, represent $F$.]

By the Paley--Wiener theorem, $F[F](x)=0$ for $|x|>\pi$, hence $F[F]\in L_{1}(-\infty,\infty)\cap L_{2}(-\infty,\infty)$. (Note that $F[F](\alpha)=F[F](-\alpha)$.) Since $(\varphi_{n})$ is complete in $L_{2}(-\pi,\pi)$, given $\varepsilon>0$, there exists a polynomial in the $\varphi_{n}$'s, with complex coefficients $\alpha_{n}$, such that
\[
\Bigl\|\check F[F]-\sum_{n=-N}^{N}\alpha_{n}\varphi_{n}\Bigr\|_{2}<\varepsilon,
\]
since $\varphi_{n}$ vanishes outside $(-\pi,\pi)$ for every $n\ge 0$. By Plancherel's theorem we have
\[
\Bigl\|F-\sum_{n=-N}^{N}\alpha_{n}\check F[\varphi_{n}]\Bigr\|_{2}<\varepsilon.
\]
This is equivalent to saying that
\begin{equation}\tag{9.4}\label{eq:II-9.4}
\|F\|_{2}^{2}=\sum_{n=-\infty}^{\infty}|a_{n}|^{2}<\infty.
\end{equation}
(See \eqref{eq:II-8.17} and the Remark after the Riesz--Fischer theorem.)

The series
\[
\sum_{n=-\infty}^{\infty}\frac{\{\sin(\pi x)\}^{2}}{(x-n)^{2}}
\]
is uniformly convergent on $(-\infty,\infty)$. By Schwarz's inequality, and \eqref{eq:II-9.4}, it follows that the series
\[
\sum_{n=-\infty}^{\infty}a_{n}(-1)^{n}\frac{\sin\pi x}{\pi(x-n)}
\]
converges uniformly on $(-\infty,\infty)$. Since $F\in E^{\pi}$ by hypothesis, $F$ is continuous, hence
\begin{equation}\tag{9.5}\label{eq:II-9.5}
F(x)=\sum_{n=-\infty}^{\infty}a_{n}(-1)^{n}\frac{\sin\pi x}{\pi(x-n)}
=\frac{\sin\pi x}{\pi}\sum_{n=-\infty}^{\infty}\frac{a_{n}(-1)^{n}}{x-n},
\end{equation}
for $-\infty<x<\infty$.

If we put $x=n$, we get
\[
a_{n}=F(n),\qquad n=0,\pm1,\pm2,\dots.
\]
The series on the right-hand side of \eqref{eq:II-9.5} converges absolutely and uniformly also with a complex $z$ in place of the real $x$, provided that $-\infty<-A<\operatorname{Im}z<A<+\infty$, hence represents an entire function of $z$, which coincides with $F(z)$ because of \eqref{eq:II-9.5}, and \eqref{eq:II-9.2} follows.
\end{proof}

Theorem~\ref{thm:II-11} can be used to prove an interpolation formula where the assumption that $F(x)\in L_{2}(-\infty<x<\infty)$ is replaced by the boundedness of $F(x)$.

\begin{theorem}\label{thm:II-12}
If $F(z)\in E^{\pi}$, $z=x+iy$, and $F(x)$ is bounded for $-\infty<x<\infty$, then we have
\begin{equation}\tag{9.6}\label{eq:II-9.6}
F(z)=\frac{\sin\pi z}{\pi}\Bigl\{F'(0)+\frac{F(0)}{\pi z}
+\sum_{\substack{n=-\infty\\ n\neq 0}}^{\infty}(-1)^{n}F(n)\Bigl(\frac{1}{z-n}+\frac{1}{n}\Bigr)\Bigr\},
\end{equation}
where $F'$ denotes the derivative of $F$.
\end{theorem}
\begin{proof}
If we set $G(z)=\frac{F(z)-F(0)}{z}$, then $G(z)\in E^{\pi}$, and $G(x)\in L_{1}(-\infty<x<\infty)$. Theorem~\ref{thm:II-11} then gives
\[
G(z)=\frac{\sin\pi z}{\pi}\sum_{n=-\infty}^{\infty}\frac{(-1)^{n}G(n)}{z-n},
\]
so that
\begin{equation}\tag{9.7}\label{eq:II-9.7}
F(z)-F(0)=\frac{z\sin\pi z}{\pi}\sum_{\substack{n=-\infty\\ n\neq 0}}^{\infty}\frac{(-1)^{n}(F(n)-F(0))}{n(z-n)}+F'(0)\frac{\sin\pi z}{\pi},
\end{equation}
since $G(0)=F'(0)$. We have, however,
\[
\frac{\pi}{\sin\pi z}=\frac{1}{z}+\sum_{\substack{n=-\infty\\ n\neq 0}}^{\infty}(-1)^{n}\Bigl\{\frac{1}{z-n}+\frac{1}{n}\Bigr\}
=\frac{1}{z}+z\sum_{\substack{n=-\infty\\ n\neq 0}}^{\infty}\frac{(-1)^{n}}{n(z-n)},
\]
so that
\[
F(0)=F(0)\frac{\sin\pi z}{\pi}\frac{\pi}{\sin\pi z}
=F(0)\frac{\sin\pi z}{\pi}\Bigl(\frac{1}{z}+z\sum_{\substack{n=-\infty\\ n\neq 0}}^{\infty}\frac{(-1)^{n}}{n(z-n)}\Bigr),
\]
and if we use this in \eqref{eq:II-9.7} we obtain \eqref{eq:II-9.6}.
\end{proof}

\subsection*{Two inequalities due to S.\ Bernstein}
\noindent The interpolation formula obtained in Theorem~\ref{thm:II-12} can be used to prove two important inequalities originally proved by S.\ Bernstein.

\begin{theorem}[S.\ Bernstein]\label{thm:II-13}
If $F(z)\in E^{\sigma}$, $\sigma>0$, $z=x+iy$, and $F(x)$ is bounded for $-\infty<x<\infty$, with $M=\sup|F(x)|$, then we have
\begin{equation}\tag{10.1}\label{eq:II-10.1}
|F'(x)|\le\sigma M,\qquad-\infty<x<\infty,
\end{equation}
the dash denoting the derivative. Equality occurs in \eqref{eq:II-10.1} if and only if
\begin{equation}\tag{10.2}\label{eq:II-10.2}
F(z)=a\,e^{i\sigma z}+b\,e^{-i\sigma z},
\end{equation}
where $a$ and $b$ are complex numbers.
\end{theorem}
\begin{proof}
It is sufficient to prove the theorem for $\sigma=\pi$, for we can otherwise study $F(\frac{z\pi}{\sigma})$.

By taking $z=x$ in \eqref{eq:II-9.6}, and differentiating once, we get
\begin{equation}\tag{10.3}\label{eq:II-10.3}
F'(x)=\cos\pi x\,F_{1}(x)+\frac{\sin\pi x}{\pi}\sum_{n=-\infty}^{\infty}\frac{(-1)^{n-1}F(n)}{(x-n)^{2}},
\end{equation}
where
\[
F_{1}(x)=F'(0)+\frac{F(0)}{x}
+\sum_{\substack{n=-\infty\\ n\neq 0}}^{\infty}(-1)^{n}F(n)\Bigl\{\frac{1}{x-n}+\frac{1}{n}\Bigr\},
\]
the differentiation being justified by the fact that the differentiated series converges uniformly. On setting $x=\frac{1}{2}$, we get
\begin{equation}\tag{10.4}\label{eq:II-10.4}
F'\Bigl(\frac{1}{2}\Bigr)=\frac{4}{\pi}\sum_{n=-\infty}^{\infty}\frac{(-1)^{n-1}F(n)}{(2n-1)^{2}},
\end{equation}
so that
\begin{equation}\tag{10.5}\label{eq:II-10.5}
\Bigl|F'\Bigl(\frac{1}{2}\Bigr)\Bigr|\le\frac{4}{\pi}M\sum_{n=-\infty}^{\infty}\frac{1}{(2n-1)^{2}}
=\frac{4M}{\pi}\frac{\pi^{2}}{4}=M\pi.
\end{equation}
For any real $x_{0}$, consider
\begin{equation}\tag{10.6}\label{eq:II-10.6}
G(z)=F\Bigl(x_{0}+z-\frac{1}{2}\Bigr).
\end{equation}
We have: $G\in E^{\pi}$, $|G(x)|\le M$, $G'(\frac{1}{2})=F'(x_{0})$, $G(n)=F(x_{0}+n-\frac{1}{2})$, and \eqref{eq:II-10.5} implies that
\[
|F'(x_{0})|=|G'(\tfrac{1}{2})|\le\pi M,
\]
which proves the first part of the theorem.

To prove the second part, we apply formula \eqref{eq:II-10.4} to the function $G$ in \eqref{eq:II-10.6}, and obtain
\[
\begin{aligned}
&G'  \Bigl(\frac{1}{2}\Bigr) = F'(x_{0})\\
&= \frac{4}{\pi}\sum_{n=-\infty}^{\infty}\frac{(-1)^{n-1}G(n)}{(2n-1)^{2}} = \frac{4}{\pi}\sum_{n=-\infty}^{\infty}\frac{(-1)^{n-1}F(x_{0}+n-\frac{1}{2})}{(2n-1)^{2}}.
\end{aligned}
\]
On replacing $x_{0}$ by $x$, and $n$ by $n+1$, we get
\begin{equation}\tag{10.7}\label{eq:II-10.7}
F'(x)=\frac{4}{\pi}\sum_{n=-\infty}^{\infty}\frac{(-1)^{n}F(x+n+\frac{1}{2})}{(2n+1)^{2}}.
\end{equation}
Let us suppose that for $x=x_{1}$ we have the equality
\[
|F'(x_{1})|=\pi M,\qquad\text{or}\qquad F'(x_{1})=\pi M e^{i\alpha},\quad\alpha\text{ real}.
\]
Formula \eqref{eq:II-10.7} then gives
\begin{equation}\tag{10.8}\label{eq:II-10.8}
\frac{4}{\pi}\sum_{n=-\infty}^{\infty}\frac{(-1)^{n}F(x_{1}+n+\frac{1}{2})-\pi M e^{i\alpha}}{(2n+1)^{2}}=0,
\end{equation}
the series on the left-hand side being absolutely convergent, since $F$ is bounded on the real axis. Since $|F(x)|\le M$, we have $-i\alpha\overline{F(x_{1}+n+\frac{1}{2})}\ge 0$, and if
\[
(-1)^{n}F(x_{1}+n+\tfrac{1}{2})=\tfrac{1}{2}M e^{i\alpha}>0,
\]
hence \eqref{eq:II-10.8} gives
\[
(-1)^{n}F(x_{1}+n+\tfrac{1}{2})=M e^{i\alpha},\qquad\text{for }n=0,\pm1,\pm2,\dots.
\]
If we set $H(z)=F(x_{1}+z+\frac{1}{2})$, then $H(0)=F(x_{1}+\frac{1}{2})=M e^{i\alpha}$, and $H(n)=F(x_{1}+n+\frac{1}{2})$, so that
\begin{equation}\tag{10.9}\label{eq:II-10.9}
(-1)^{n}H(n)=H(0).
\end{equation}
If we apply Theorem~\ref{thm:II-12} to the function $H(z)$, we get
\[
H(z)=\frac{\sin\pi z}{\pi}\Bigl\{H'(0)+\frac{H(0)}{z}
+\sum_{n=-\infty}^{\infty}H(0)\Bigl(\frac{1}{z-n}+\frac{1}{n}\Bigr)\Bigr\},
\]
because of \eqref{eq:II-10.9}. We have, however,
\[
\pi\cot\pi z=\frac{1}{z}+\sum_{\substack{n=-\infty\\ n\neq 0}}^{\infty}\Bigl(\frac{1}{z-n}+\frac{1}{n}\Bigr).
\]
Hence
\[
\begin{aligned}
H(z) &= \frac{\sin\pi z}{\pi}\bigl[H'(0)+H(0)\cdot\pi\cot\pi z\bigr]\\
&= A\cos\pi z+B\sin\pi z,\quad\text{say},\\
&= A_{1}e^{\pi z i}+B_{1}e^{-\pi z i}.
\end{aligned}
\]
Since $F(z)=H(z-x_{1}-\frac{1}{2})$, it is proved that if equality occurs in \eqref{eq:II-10.1}, then \eqref{eq:II-10.2} holds.

We have to show that if \eqref{eq:II-10.2} holds, then equality occurs in \eqref{eq:II-10.1} for some $x$, $-\infty<x<\infty$.

If $F(x)=a\,e^{i\sigma x}+b\,e^{-i\sigma x}$, and $a=|a|e^{i\alpha}$, $b=|b|e^{i\beta}$, $\sigma>0$, then
\[
F(x)=e^{i(\alpha+\sigma x)}\bigl(|a|+|b|\bigr)e^{i(\beta-\alpha)-2i\sigma x},
\]
so that for $x=x_{1}=\frac{(\beta-\alpha)-\pi}{2\sigma}$, we have $|F'(x_{1})|=\sigma(|a|+|b|)=\sigma\max_{x}|F(x)|$.
\end{proof}

\begin{remark}
By letting $\sigma\downarrow 0$ in \eqref{eq:II-10.1}, we deduce that the only functions $F\in E^{0}$, which are bounded on the real axis, are constants.

On the other hand, it follows from a theorem of Siegel that if $F\in E^{0}$, and $F(x)\in L_{1}(-\infty<x<\infty)$, then $F\equiv 0$.
\end{remark}

\begin{theorem}[S.\ Bernstein]\label{thm:II-14}
Let $T(x)=\sum_{k=-n}^{n}c_{k}e^{ikx}$, where $x$ is real and $c_{k}$ complex. Then the inequality
\[
|T(x)|\le M
\]
implies that
\[
|T'(x)|\le n M,
\]
and this inequality becomes an equality if and only if
\[
T(x)=a\,e^{inx}+b\,e^{-inx},
\]
where $a$ and $b$ are complex numbers.
\end{theorem}
\begin{proof}
We have only to note that for complex $z$, $T(z)$ is an entire function, with $T\in E^{n}$, while $T(x)$ is bounded, and apply Theorem~\ref{thm:II-13}.
\end{proof}

\setcounter{section}{10}
\section{Several variables}
\label{sec:II-11}

The proof of Plancherel's theorem in several variables is not essentially different from that in one variable. If $E_{k}$ denotes the real Euclidean space of $k$ dimensions, let $\mathcal S$ denote the Schwartz space of infinitely differentiable functions on $E_{k}$, such that for any $\alpha=(\alpha_{1},\dots,\alpha_{k})$, $\beta=(\beta_{1},\dots,\beta_{k})$, where $\alpha_{1},\alpha_{2},\dots,\alpha_{k}$ and $\beta_{1},\beta_{2},\dots,\beta_{k}$ are non-negative integers,
\[
\sup_{x\in E_{k}}\bigl|x^{\alpha}(D^{\beta}f)(x)\bigr|<\infty,
\]
where $x^{\alpha}=x_{1}^{\alpha_{1}}\cdots x_{k}^{\alpha_{k}}$, and $D^{\beta}f=(\frac{\partial}{\partial x_{1}})^{\beta_{1}}\cdots(\frac{\partial}{\partial x_{k}})^{\beta_{k}}$. Then $\mathcal S$ is a dense subspace of $L_{p}(E_{k})$, $1\le p<\infty$. For $f\in\mathcal S$ we define the Fourier transform $F[f]$ by the relation
\[
F[f](x)=(2\pi)^{-k/2}\int_{E_{k}}f(t)\,e^{i\langle x,t\rangle}\,dt,\qquad x\in E_{k}.
\]
Then $F[f]\in\mathcal S$, and $f\mapsto F[f]$ is a one-to-one mapping of $\mathcal S$ onto $\mathcal S$ (as in the case of one variable, cf.\ \S\ref{sec:II-2}). If $f,g\in\mathcal S$, then $f*g\in\mathcal S$. We further have
\[
\int_{E_{k}}f(x)\,F[g](x)\,dx
=\int_{E_{k}}F[f](y)\,g(y)\,dy,
\]
and
\[
(f,g)=\bigl(F[f],F[g]\bigr).
\]
If $f\in L_{2}(E_{k})$, there exists a sequence $(f_{n})$ of functions belonging to $\mathcal S$, such that $\|f-f_{n}\|_{2}\to 0$, as $n\to\infty$, and $\|F[f_{n}]-F[f_{m}]\|_{2}=\|f_{m}-f_{n}\|_{2}$, which implies that $\|F[f_{m}]-F[f_{n}]\|_{2}\to 0$, as $m,n\to\infty$. Since $L_{2}(E_{k})$ is complete, there exists $g\in L_{2}(E_{k})$, such that $\|g-F[f_{m}]\|_{2}\to 0$, as $m\to\infty$. And $\|g\|_{2}=\lim_{m\to\infty}\|F[f_{m}]\|_{2}=\lim_{m\to\infty}\|f_{m}\|_{2}=\|f\|_{2}$. We define $g$ to be the Fourier transform of $f\in L_{2}(E_{k})$. It is independent of the approximating sequence, and is defined almost everywhere. We denote it by $F[f]$.

We similarly define
\[
\check F[f](x)=(2\pi)^{-k/2}\int_{E_{k}}f(t)e^{-i\langle x,t\rangle}\,dt,
\]
and extend the definition to all of $L_{2}(E_{k})$. It follows, as in \eqref{eq:II-2.11b}, that
\[
\check F\bigl[F[f]\bigr]=f=F\bigl[\check F[f]\bigr],
\]
so that $f\mapsto F[f]$ is a linear mapping of $L_{2}(E_{k})$ onto itself; it is also isometric. The proof that the definitions of the Fourier transform on $L_{2}(E_{k})$, and on $L_{1}(E_{k})$ coincide on $L_{1}(E_{k})\cap L_{2}(E_{k})$ follows as in the case of $E_{1}$.
